\documentclass[UTF8,fontset=windows,a4paper,11pt]{ctexart}

% ---- page geometry ----
\usepackage[margin=2.5cm]{geometry}

% ---- math ----
\usepackage{amsmath,amssymb}

% ---- fonts: CJK via ctex windows fontset (宋体 SimSun / bold 黑体 SimHei /
%      italic 楷体 KaiTi); Latin text stays Latin Modern (good Latin glyphs) ----
\setmonofont{Consolas}

% ---- graphics / captioned figure floats ----
\usepackage{graphicx}
% Paper carries manual, non-sequential labels (图 1, 图 4-1, 图 7-6) inside the
% caption text, so strip the kernel "Figure N:" prefix. We redefine \@makecaption
% directly rather than load the caption package, which conflicts with pandoc's
% caption-less longtables ("No counter '' defined").
\makeatletter
\def\fps@figure{htbp}
\long\def\@makecaption#1#2{%
  \vskip\abovecaptionskip
  \begingroup\small
    \sbox\@tempboxa{#2}%
    \ifdim\wd\@tempboxa>\hsize #2\par
    \else \hb@xt@\hsize{\hfil\box\@tempboxa\hfil}\fi
  \endgroup
  \vskip\belowcaptionskip}
\makeatother

% ---- tables (pandoc longtable + booktabs) ----
\usepackage{longtable,booktabs,array,calc}
\usepackage{etoolbox}
\makeatletter
\patchcmd\longtable{\par}{\if@noskipsec\mbox{}\fi\par}{}{}
\makeatother
\IfFileExists{footnotehyper.sty}{\usepackage{footnotehyper}}{\usepackage{footnote}}
\makesavenoteenv{longtable}
\providecommand{\real}[1]{#1}

% ---- pandoc helper macros (replicated, since we use a custom template) ----
\makeatletter
\newsavebox\pandoc@box
\newcommand*\pandocbounded[1]{%
  \sbox\pandoc@box{#1}%
  \Gscale@div\@tempa{\textheight}{\dimexpr\ht\pandoc@box+\dp\pandoc@box\relax}%
  \Gscale@div\@tempb{\linewidth}{\wd\pandoc@box}%
  \ifdim\@tempb\p@<\@tempa\p@\let\@tempa\@tempb\fi%
  \ifdim\@tempa\p@<\p@\scalebox{\@tempa}{\usebox\pandoc@box}%
  \else\usebox{\pandoc@box}%
  \fi%
}
\makeatother
\providecommand{\tightlist}{\setlength{\itemsep}{0pt}\setlength{\parskip}{0pt}}
\setlength{\emergencystretch}{3em}

% ---- paragraph style ----
\setlength{\parindent}{0pt}
\setlength{\parskip}{6pt plus 2pt minus 1pt}

% ---- section numbering off (paper carries manual 第X章 / X.Y numbers) ----
\setcounter{secnumdepth}{-\maxdimen}
\setcounter{tocdepth}{2}

% ---- glyph mapping: render math symbols in math mode regardless of font ----
\usepackage{newunicodechar}
% circled digits -> SimSun (Latin font lmroman lacks U+2460..; \char avoids active-char recursion)
\newfontfamily\cjkcircled{SimSun}[AutoFakeBold=3]
\newunicodechar{①}{{\cjkcircled\char"2460}}
\newunicodechar{②}{{\cjkcircled\char"2461}}
\newunicodechar{③}{{\cjkcircled\char"2462}}
% math symbols -> math mode (font-independent, correct typography)
\newunicodechar{−}{\ensuremath{-}}
\newunicodechar{≈}{\ensuremath{\approx}}
\newunicodechar{≤}{\ensuremath{\leq}}
\newunicodechar{≥}{\ensuremath{\geq}}
\newunicodechar{≠}{\ensuremath{\neq}}
\newunicodechar{∈}{\ensuremath{\in}}
\newunicodechar{ρ}{\ensuremath{\rho}}
\newunicodechar{Δ}{\ensuremath{\Delta}}
\newunicodechar{Ω}{\ensuremath{\Omega}}
\newunicodechar{·}{\ensuremath{\cdot}}
\newunicodechar{×}{\ensuremath{\times}}
\newunicodechar{±}{\ensuremath{\pm}}
\newunicodechar{→}{\ensuremath{\rightarrow}}
\newunicodechar{′}{\ensuremath{\prime}}
\newunicodechar{²}{\textsuperscript{2}}
\newunicodechar{₁}{\textsubscript{1}}
\newunicodechar{₂}{\textsubscript{2}}

% ---- hyperref ----
\usepackage[unicode]{hyperref}
\hypersetup{
  pdftitle={DAG-BFT
共识中认证机制的结构性作用：去认证设计的隐性代价与性能边界},
  colorlinks=true, linkcolor={black}, urlcolor={blue!55!black}, citecolor={black},
  pdfborder={0 0 0}}

\title{DAG-BFT
共识中认证机制的结构性作用：去认证设计的隐性代价与性能边界}
\author{董孟哲\\523030910105}
\date{2026.06}

\begin{document}
\maketitle
\tableofcontents
\bigskip

\section{摘要}\label{ux6458ux8981}

有向无环图(DAG)型拜占庭容错(BFT)共识近年从 certified 转向
uncertified,以移除每区块认证换取更低延迟;这一优势多被视为结构性的------少一轮可靠广播(RBC)即少一轮延迟。本文质疑此判断。现有
certified 与 uncertified
的比较混淆变量未受控:二者除认证外尚有十余处设计差异,将延迟差直接归因于去认证并不可靠。为此本文提出一个以安全性与活性不变量为硬边界的受控分析框架,隔离认证强度这一单一变量,并论证:移除认证不是可分离的性能优化,而是一次结构性替换------认证承载的三个不变量(非等价性、数据可用性、因果历史完整性)并未消失,只是被重新分配到其他组件。关键的不对称在于:逻辑型不变量可被低成本平移,物理型的数据可用性只能被弱化并补偿;隐性代价因而集中于可用性,且仅在故障时爆发,是一笔或有负债而非固定开销。本文进一步在异步模型下论证非等价性必须被
enforce,从而说明 uncertified
的延迟优势是条件性的。真正决定协议行为的深层设计轴,是可用性的
enforcement 机制与协调修复触发条件,而非 certified/uncertified
这一浅层标签。该结论的或有负债侧已由真实 Sui n=7/n=4
部署的受控故障注入(290+
次重复)初步支持:故障代价确为非线性、由故障强度驱动的或有负债(详见
§7.6.1)。

\textbf{关键词}:区块链共识;DAG-based BFT;Reliable Broadcast;uncertified
DAG;数据可用性;协议不变量;延迟分析

\section{Abstract}\label{abstract}

DAG-based Byzantine Fault Tolerant (BFT) consensus has recently shifted
from certified to uncertified DAGs, trading per-block certification for
lower commit latency---an advantage widely treated as structural: one
fewer round of Reliable Broadcast (RBC) means one fewer round of
latency. This paper challenges that judgment. Existing
certified-versus-uncertified comparisons are confounded: the two
families differ along more than a dozen design dimensions, so
attributing the latency gap to de-certification alone is unsafe. We
therefore propose a controlled analytical framework that draws module
boundaries along safety and liveness invariants, isolating certification
strength as a single variable, and argue that removing certification is
not a separable optimization but a structural substitution: the three
invariants it carries---non-equivocation, data availability, and
causal-history integrity---are reassigned rather than eliminated. The
key asymmetry is that logical invariants relocate cheaply, whereas
physical data availability can only be weakened and compensated, so the
hidden cost concentrates on availability as a contingent liability that
materializes only under faults. We further argue, in the asynchronous
model, that non-equivocation must be enforced, making the latency
advantage conditional. The genuine design axis is the
availability-enforcement mechanism and its reconciliation trigger, not
the surface certified/uncertified label. This contingent-liability side
is preliminarily supported by controlled fault-injection on a real Sui
n=7/n=4 deployment (290+ repetitions): the fault-induced latency cost is
a non-linear, fault-intensity-driven contingent liability.

\textbf{Keywords}: blockchain consensus; DAG-based BFT; Reliable
Broadcast; uncertified DAG; data availability; protocol invariants;
latency analysis

\section{第 1 章 引言}\label{ux7b2c-1-ux7ae0-ux5f15ux8a00}

\subsection{1.1 研究背景}\label{ux7814ux7a76ux80ccux666f}

区块链系统的吞吐量与延迟长期受制于其共识层。经典的拜占庭容错共识协议,如
PBFT {[}10{]} 与 HotStuff {[}11{]},采用 leader-based 结构:由单一 leader
打包区块并驱动多轮投票。这一结构在工程上简洁,但 leader
同时是带宽瓶颈与性能单点------所有交易数据需经由 leader 分发,leader
的速度直接决定全网速度
{[}33{]}。当验证者规模扩大、交易负载升高时,这一瓶颈愈发突出。

DAG-based BFT
共识为打破该瓶颈提供了一条新路径。其核心思想是将``数据传播''与``共识排序''解耦
{[}31{]},
{[}32{]}:每个验证者并行地广播自己的区块,每个区块引用一批更早的区块,这些引用关系自然织成一张有向无环图;随后所有验证者对同一张本地
DAG 运行一个确定性的排序规则,无需额外通信即可得到一致的全序 {[}31{]}。自
DAG-Rider {[}31{]} 在 2021 年给出异步 DAG
共识的理论基础以来,Narwhal/Tusk {[}32{]}、Bullshark {[}33{]}
等协议将其工程化落地,并被 Sui、Aptos 等公链采用 {[}38{]}, {[}46{]}。

然而,DAG-based BFT 也带来了新的代价:延迟。由于 DAG
逐层推进,而每一层的区块此前需经由可靠广播(RBC)认证,提交一个区块往往需要等待多轮
RBC。已有测量表明,Tusk 的良好情形延迟约为 7 个通信步、最坏情形约 21
步,DAG-Rider 的最优延迟约 12 步,均显著高于 PBFT 的 3 步 {[}35{]},
{[}36{]}。这一``延迟焦虑''催生了一系列优化工作:Shoal 与 Shoal++
通过流水线化降低锚点延迟 {[}37{]}, {[}38{]},Sailfish 让每一轮都拥有
leader 顶点 {[}39{]}。

\subsection{1.2
一个被默认的判断}\label{ux4e00ux4e2aux88abux9ed8ux8ba4ux7684ux5224ux65ad}

在上述优化中,最激进的一条路线是放弃认证本身。Cordial Miners {[}42{]}
弃用 RBC,改以 blocklace 结构实现共识,将良好情形延迟由 4 步降至 3
步;Mysticeti {[}43{]} 进一步提出 uncertified
DAG,移除对区块的显式认证,据称达到 3 个消息轮次的延迟下界,在广域网下实现
0.5 秒提交延迟与超过 20 万 TPS 的吞吐;Mahi-Mahi {[}44{]} 亦在
uncertified DAG 上构建低延迟异步共识。

这些工作传递出一个被广泛接受的判断:uncertified DAG
主要通过去掉认证机制来``免费''降低延迟------少一轮
RBC,即少一轮延迟。这一判断在直觉上极具说服力,也几乎成为该方向的共识性叙述。

本文质疑这一判断,质疑的起点是一个方法论问题。本文记之为
\textbf{SubRQ0:certified 与 uncertified DAG 的比较对象是否公平?}
现有比较通常将某个 certified 协议(如 Bullshark {[}33{]})与某个
uncertified 协议(如 Mysticeti
{[}43{]})并置,比较其报告的延迟数字。但这两类协议除``是否认证''之外,在提交规则、leader
调度、流水线化程度、网络模型假设等十余个维度上均不相同。当 uncertified
协议表现出更低延迟时,本文无法确定这一差异究竟来自``去认证'',还是来自其他被一并改动的设计选择。换言之,现有比较是一种混淆变量未受控的比较,其结论的内部效度存疑。

\subsection{1.3
本文论点与贡献}\label{ux672cux6587ux8bbaux70b9ux4e0eux8d21ux732e}

本文的中心论点是:\textbf{移除认证并非一项可分离的性能优化,而是一次结构性替换。}
认证所承载的安全性与活性不变量并未因移除而消失,而是被重新分配至协议的其他组件;因此
uncertified 相对 certified 的延迟优势是条件性的,而非结构性的。

本文的贡献分为递进的三层:

\textbf{贡献 A(benchmark
批判)}:本文论证,现有以延迟数字判定协议优劣的比较,实质上是对一个高维性能函数在少数维度上的投影,且现有评测流程系统性地偏向
happy-path 附近的采样;certified 与 uncertified
的相对优劣会随故障模式、网络非对称与负载偏斜而改变。

\textbf{贡献 B(fundamentality 批判)}:本文论证,uncertified
的良好情形延迟优势并非结构性的,而会在恢复频率升高、部分同步退化、拥塞反馈等条件下被补偿机制的摊销成本反向吞噬。这一或有负债的存在与非线性形态已获一项真实部署实证的初步支持(§7.6.1)。

\textbf{贡献 C(轴批判,主贡献)}:本文论证,certified/uncertified
并非描述协议的合适坐标轴;真正决定协议行为的深层设计轴是
\textbf{availability-enforcement
机制}(数据可用性以何种方式、在何时被强制)\textbf{与 reconciliation
trigger 条件}(数据缺失时的协调修复以何种条件触发)。certified 与
uncertified 只是这一深层轴上的两个投影点。

围绕这一论点,本文回答以下研究问题:

\begin{quote}
\textbf{主研究问题(RQ)}:在何种网络与敌手条件下,uncertified DAG-BFT 相对
certified 设计的``低延迟优势''才真正成立?

\begin{itemize}
\tightlist
\item
  \textbf{SubRQ1}:认证机制为协议买到了哪些理论保证?
\item
  \textbf{SubRQ2}:去掉认证后,这些保证被转移或重建到哪里?隐性代价以何种形式出现?
\item
  \textbf{SubRQ3}:延迟优势在哪些条件下成立、被抵消、甚至反转?
\end{itemize}
\end{quote}

须说明本文给出的是一个条件化结论,而条件化结论存在被断章取义的风险:读者可能抽取``uncertified
没问题''这一句而略去其前提条件,据此为真实部署中削减认证的决策背书。本文在此明确声明:本文的任何结论都不应脱离其网络模型与敌手模型前提被引用。值得说明的是,本文的主贡献
C 在表述形式上具有一定的抗断章性------它陈述的是一个设计轴而非一个``X
优于 Y''的比较句,引用者若要引用该轴,必须连同轴的定义一并引用。

\subsection{1.4 论文结构}\label{ux8bbaux6587ux7ed3ux6784}

第 2 章梳理区块链共识与 DAG 共识的谱系并界定研究空白;第 3
章提出以不变量为边界的受控分析方法论;第 4 章回答 SubRQ1;第 5 章回答
SubRQ2,为全文核心;第 6 章给出非等价性必要性的理论论证;第 7 章回答
SubRQ3,并以 §7.6.1 的真实部署实证为其或有负债结论提供经验锚点;第 8
章提出主贡献 C 并讨论其对排序公平性的延伸;第 9 章讨论局限与威胁有效性;第
10 章总结。

\section{第 2 章
背景与相关工作}\label{ux7b2c-2-ux7ae0-ux80ccux666fux4e0eux76f8ux5173ux5de5ux4f5c}

\subsection{2.1
区块链共识的两大谱系}\label{ux533aux5757ux94feux5171ux8bc6ux7684ux4e24ux5927ux8c31ux7cfb}

区块链共识协议大致可分为两大谱系。第一条谱系源于 Nakamoto 的 Bitcoin
设计 {[}22{]},以工作量证明与最长链规则达成概率性一致,其确认延迟以分钟乃至小时计 {[}23{]}, {[}24{]}。第二条谱系源于经典分布式系统的状态机复制(SMR) {[}4{]},追求确定性最终性与秒级乃至亚秒级延迟,本文研究的 DAG-based BFT 即属于此。值得注意的是,``DAG''一词在两条谱系中含义不同:在 Nakamoto 谱系中它指 SPECTRE/PHANTOM 式的区块 DAG,在本文语境中它特指 BFT 共识的通信结构。

\subsection{2.2
拜占庭容错共识的理论约束与经典协议}\label{ux62dcux5360ux5eadux5bb9ux9519ux5171ux8bc6ux7684ux7406ux8bbaux7ea6ux675fux4e0eux7ecfux5178ux534fux8bae}

拜占庭容错问题最早由 Lamport 等人形式化 {[}1{]}:在 n 个进程中存在至多 f
个表现任意(包括恶意)的故障进程时,正确进程仍需就某一值达成一致。共识的两个核心性质是安全性(safety,正确进程不会产生分歧的决定)与活性(liveness,决定最终会做出)。Fischer
等人的 FLP
不可能性结果指出,在完全异步且允许一个进程崩溃的系统中,不存在同时保证三者的确定性共识算法
{[}2{]}。这一结果迫使后续工作要么引入时间假设,要么引入随机性。

Dwork 等人提出的部分同步模型 {[}3{]}
成为最具影响力的折中:系统在某个未知的全局稳定时刻(GST)之后表现同步。PBFT
{[}10{]} 是部分同步模型下首个实用的 BFT 协议,HotStuff {[}11{]}
将其通信复杂度降为线性。与部分同步路线并行的是异步 BFT {[}7{]}, {[}17{]},其代价通常是更高的常数因子与更复杂的协议逻辑。门限签名等密码学基元 {[}8{]}------为这些协议提供了把 O(n)
个签名压缩为常数大小证书的能力,是认证机制工程实现的基础。

\subsection{2.3 DAG-based BFT
的核心思想}\label{dag-based-bft-ux7684ux6838ux5fc3ux601dux60f3}

leader-based
协议的根本瓶颈在于:数据分发与共识排序由同一条关键路径承担,而该路径以
leader 为中心。DAG-based BFT 的突破在于将二者彻底解耦 {[}31{]},
{[}32{]}。

在 DAG
共识中,系统按轮推进。每一轮中,每个验证者构造一个区块,该区块包含交易数据以及对上一轮中至少
2f+1
个区块的引用(称为边)。所有验证者并行广播各自的区块,这些区块与引用关系共同构成一张有向无环图。关键之处在于:由于所有正确验证者最终会看到同一张(或足够一致的)DAG,它们可以各自在本地运行一个确定性的排序规则,把图``翻译''成一个一致的全序,而这一翻译过程几乎不需要额外通信
{[}31{]}。数据传播因此可以全速并行,排序的通信开销被压至最低。Hashgraph
{[}61{]} 是这一思路的一个早期工业实例,它用``gossip about
gossip''构造类似结构并以虚拟投票完成排序。

\subsection{2.5 certified DAG
的谱系}\label{certified-dag-ux7684ux8c31ux7cfb}

DAG-Rider {[}31{]}
是这一思路的理论奠基:下层让每个进程通过可靠广播(RBC)广播其提案并构建 DAG,上层以零额外通信完成全序。RBC 保证一个区块一旦被某个正确进程接受,就会被所有正确进程接受------本文称经由 RBC 广播的区块为``认证的''(certified)区块。Narwhal/Tusk {[}32{]} 将该思路工程化;Bullshark {[}33{]} 在部分同步模型下提供低延迟快速路径。其后 Aleph {[}34{]}、Shoal/Shoal++ {[}37{]}, {[}38{]}、Sailfish {[}39{]} 等在广播原语与延迟上持续优化 {[}35{]}, {[}36{]}, {[}40{]}, {[}41{]}。这一谱系的共同特征是:它们都保留了对区块的某种认证。

\subsection{2.6 uncertified DAG
的转向}\label{uncertified-dag-ux7684ux8f6cux5411}

认证带来保证,也带来延迟:Bracha 的经典 RBC 需要两个阶段、O(n²) 条消息 {[}5{]}, {[}9{]}。为削减这一开销,一条新路线选择放弃认证。Cordial Miners {[}42{]} 弃用 RBC,将良好情形延迟由 4 步降至 3 步;Mysticeti {[}43{]} 提出 uncertified DAG,据称达到 3 个消息轮次的延迟下界;Mahi-Mahi {[}44{]} 在 uncertified 结构化 DAG 上构建异步共识。值得注意的是,uncertified DAG 并非没有代价:Adelie {[}45{]} 明确指出 Mysticeti 的 uncertified DAG 引入了 certified DAG 中并不存在的新拜占庭攻击向量,并为此设计了专门的检测与防护机制。这一观察是本文论点的重要经验支撑。

\subsection{2.7 研究空白}\label{ux7814ux7a76ux7a7aux767d}

综观上述谱系,已有工作的主流范式是``提出一个新协议并报告其 benchmark''。即便是系统化综述 {[}57{]}--{[}60{]} 也主要在分类与梳理协议,而尚无工作把 certified/uncertified 之分本身当作分析对象、追问它是否遮蔽了一个更恰当的坐标并加以重构。Mysticeti {[}43{]}、Adelie {[}45{]} 等单个协议工作确实讨论过认证的取舍,但它们讨论的是``在这条轴上如何取值'',而非``这条轴本身是否是描述设计空间的恰当坐标''。本文要填补的正是后一种 reframing 的空白。

\section{第 3 章
方法论:不变量分解与受控比较框架}\label{ux7b2c-3-ux7ae0-ux65b9ux6cd5ux8bbaux4e0dux53d8ux91cfux5206ux89e3ux4e0eux53d7ux63a7ux6bd4ux8f83ux6846ux67b6}

\subsection{3.1
SubRQ0:为何``硬比两个异质系统''会出问题}\label{subrq0ux4e3aux4f55ux786cux6bd4ux4e24ux4e2aux5f02ux8d28ux7cfbux7edfux4f1aux51faux95eeux9898}

回到第 1 章提出的 SubRQ0。把 Bullshark {[}33{]} 与 Mysticeti {[}43{]}
的延迟数字并置比较,在方法论上类似于一个未设对照的实验:两个被比较对象在``是否认证''之外还相差众多设计维度,这些维度构成了混淆变量。若
uncertified
协议恰好同时采用了更激进的流水线化或更宽松的网络假设,则其延迟优势的一部分(甚至大部分)可能应归因于这些因素,而非``去认证''本身。把全部差异记在``去认证''账上,是一种归因错误。

因此,任何宣称``去认证降低了延迟''的论断,都负有一项方法论义务:给出一种能够单独隔离``认证''这一变量的比较方式。否则,SubRQ0
只是一句抱怨,而非一项研究。

\subsection{3.2
以不变量为模块硬边界的分解}\label{ux4ee5ux4e0dux53d8ux91cfux4e3aux6a21ux5757ux786cux8fb9ux754cux7684ux5206ux89e3}

要隔离单一变量,实验科学的做法是控制变量;理论分析中与之对应的做法是模块化分解------把复杂协议抽象为若干功能模块,再构造一个``仅改变认证强度''的比较维度。但模块化分解本身存在一个研究者自由度:模块的边界由分析者划定,而划界是一个带假设的选择。若边界划得不当,分析者可能把一个混淆变量恰好藏进某个模块内部,从而以一个新的主观性复活了
SubRQ0 想要消除的不公平。

本文用一个客观标准来消除这一自由度:\textbf{模块的硬边界,按协议为满足安全性与活性所必须维持的不变量来划定。}
每一个模块,对应``维持某一个不变量''的机制。这一标准之所以客观,是因为安全性与活性不变量是协议为正确性所必须维持的属性,它们由协议的正确性目标决定,而不由分析者裁量。依此划界,分解便不再是主观选择。

\subsection{3.3 受控 minimal-pair
比较}\label{ux53d7ux63a7-minimal-pair-ux6bd4ux8f83}

在不变量分解的基础上,本文采用受控 minimal-pair
比较:不把两个完全不同的系统硬比,而是在尽量相同的协议背景下,观察``去认证''这一单一改动所引起的保证迁移与代价变化。理想的
minimal pair
是同一协议骨架的``认证版''与``去认证版'',二者仅在认证强度上不同。

\subsubsection{3.3.1 语料库 C 的界定}\label{ux8bedux6599ux5e93-c-ux7684ux754cux5b9a}

本文的归纳性分析建立在一个显式界定的协议语料库 C 之上。C 的纳入标准为:(i)以 DAG 为通信结构的 BFT 共识协议,(ii)有公开发表的安全性或活性证明(或等价的形式化论证),(iii)有明确的提交规则描述。据此,C 包含以下七个协议:DAG-Rider {[}31{]}、Narwhal/Tusk {[}32{]}、Bullshark {[}33{]}、Cordial Miners {[}42{]}、Mysticeti {[}43{]}、Mahi-Mahi {[}44{]}、Adelie {[}45{]}。其中前三个为 certified DAG,后四个为 uncertified DAG 或其变体。排除标准为:无安全性证明的原型系统、未发表的工作、以及非 DAG 结构的 BFT 协议(如 PBFT、HotStuff)。本文所有``语料库 C 中的协议均\ldots\ldots''形式的断言,均限定在上述七个协议之内。

\subsection{3.4
三分判据与成本-证明双判据}\label{ux4e09ux5206ux5224ux636eux4e0eux6210ux672c-ux8bc1ux660eux53ccux5224ux636e}

本文需要一个判据来回答``某个不变量在去认证后究竟发生了什么''。一个朴素的二分法------``被重新分配''或``被消除''------是不够的:若``重新分配''没有精确判据,该判断将永远不可能被证伪(总能在协议里指认某个组件``接手了''该不变量),从而沦为同义反复。

本文采用\textbf{三分判据(C1-B)}:某个不变量在去认证后,其命运必属以下三类之一。

\begin{itemize}
\tightlist
\item
  \textbf{① 平移}:该不变量被一个不同的机制以同等强度维持。
\item
  \textbf{②
  削弱+补偿}:该不变量的保证被降级,但协议引入了补偿机制,把降级造成的损害约束在可控范围内。
\item
  \textbf{③ 消除}:该不变量的保证被真正放弃,且无补偿机制。
\end{itemize}

本文的中心主张可据此精确表述为:\textbf{去认证只产生 ① 与 ②,从不产生
③;而隐性代价集中在 ②。}

这里须对该主张的非平凡性作一个诚实的界定。``去认证从不产生
③''本身并非一个令人惊讶的发现:任何能正常运行的协议都是安全且活的,其安全性/活性不变量按定义必被某种机制维持,因此
③(无补偿的消除)对一个工作协议而言几乎必然不出现。也就是说,①/②
二分本身更应被理解为本框架的\textbf{起点},而非结论。本文真正非平凡、且并非同义反复的主张有两点:其一,\textbf{①
与 ② 的落点不是随机的,而可由第 4.5
节的逻辑型/物理型二分提前预测}------逻辑型不变量系统性地落入
①,物理型的可用性系统性地落入
②;这个``可预测的不对称''是一个有内容的经验论断,它本可以不成立(例如某个逻辑型不变量可能因协议结构而无法低成本平移)。其二,\textbf{落入
②
的可用性,其补偿成本是一笔或有负债而非固定开销}------它在良好情形为零、在故障情形集中爆发,这一成本结构正是第
7 章条件化结论的根源。后文的 ①②③ 归类服务于确立这两点,而非服务于``③
不发生''这一近乎定义的陈述。

为使 ①②③
的归类可操作、可证伪,本文辅以\textbf{成本-证明双判据(C1-C)}:一个不变量被判定为``被重新承接'',当且仅当(a)在协议的成本台账中能找到一笔可归因于维持该不变量(哪怕是其弱化版本)的非零成本,且(b)在协议的安全性或活性证明中存在一个步骤明确依赖于该不变量的某个版本。判据(a)排除``嘴上说接手、实则零成本''的虚假平移,判据(b)排除``成本存在但与正确性无关''的误判。两个判据必须同时满足,缺一不可:仅满足(a)可能是与正确性无关的工程开销,仅满足(b)可能是一个未被实现却被证明默认的理想假设。

为便于对照,下表汇总两组判据:

{\def\LTcaptype{} % do not increment counter
\begin{longtable}[]{@{}
  >{\raggedright\arraybackslash}p{(\linewidth - 4\tabcolsep) * \real{0.3333}}
  >{\raggedright\arraybackslash}p{(\linewidth - 4\tabcolsep) * \real{0.3333}}
  >{\raggedright\arraybackslash}p{(\linewidth - 4\tabcolsep) * \real{0.3333}}@{}}
\toprule\noalign{}
\begin{minipage}[b]{\linewidth}\raggedright
判据
\end{minipage} & \begin{minipage}[b]{\linewidth}\raggedright
作用
\end{minipage} & \begin{minipage}[b]{\linewidth}\raggedright
内容
\end{minipage} \\
\midrule\noalign{}
\endhead
\bottomrule\noalign{}
\endlastfoot
C1-B 三分判据 & 刻画不变量去认证后的命运 & ① 平移(等强度、机制不同)/ ②
削弱+补偿(保证降级,补偿机制约束损害)/ ③ 消除(放弃且无补偿) \\
C1-C 判据(a)成本 & 排除``零成本的虚假平移'' &
成本台账中存在一笔可归因于维持该不变量(含其弱化版本)的非零成本 \\
C1-C 判据(b)证明 & 排除``与正确性无关的开销'' &
协议安全性/活性证明中存在一步明确依赖该不变量的某个版本 \\
\end{longtable}
}

判定一个不变量``被重新承接'',须 C1-C
的(a)(b)同时满足;再据其保证强度是否降级,归入 ① 或 ②。

\subsection{3.5
方法论的自反性}\label{ux65b9ux6cd5ux8bbaux7684ux81eaux53cdux6027}

这一框架有一个值得正视的自反性后果。若沿用 minimal-pair
的设想,本文似乎需要一个真实存在的``同一协议的认证版与去认证版''。但本文将在第
5 章论证:这样的 minimal pair
在严格意义上往往无法构造------因为认证不能被孤立地拿掉。真实的
uncertified
协议从来不是``某协议减去认证'',而是``某协议减去认证、再加上一批补偿机制''。这一不可构造性并非本文方法论的失败,恰恰相反,它是本文核心发现的前奏:认证不是一个可独立拆卸的模块,而是一个结构性机制。换言之,方法论上的障碍本身,转化成了一项实质发现。

\section{第 4 章
SubRQ1:认证机制买到的三个不变量}\label{ux7b2c-4-ux7ae0-subrq1ux8ba4ux8bc1ux673aux5236ux4e70ux5230ux7684ux4e09ux4e2aux4e0dux53d8ux91cf}

\subsection{4.1
认证即可靠广播}\label{ux8ba4ux8bc1ux5373ux53efux9760ux5e7fux64ad}

在 certified DAG 中,认证由可靠广播(RBC)实现。Bracha 的经典 RBC {[}5{]}
是一个两阶段协议:发送者广播消息,各进程经过 echo 与 ready
两轮交互后才``交付''该消息;在 n≥3f+1 时,它以 O(n²)
条消息保证若干性质。Cachin 等人的著作系统地给出了这些性质的形式化定义
{[}9{]}:正确性、一致性、完整性与终止性。在工程实现中,RBC
所产生的``认证''通常物化为一个由 2f+1 个验证者签名聚合而成的证书,BLS
聚合签名 {[}8{]} 使该证书的大小不随验证者数量增长。在 DAG
共识语境下,一个区块经由 RBC
广播后即成为``认证区块'',并随之获得三项关键保证。本文把这三项保证识别为认证所维持的三个不变量。须说明这一识别的范围:本文聚焦的三个不变量------非等价性、数据可用性、因果历史完整性------分别对应
RBC 的一致性、终止性与完整性三条性质在 DAG
语境下的投影,它们是去认证分析中真正承载安全性/活性风险的维度。RBC
的标准性质集还包含一条 validity/correctness
性质(正确发送者广播的消息终将被正确进程交付);本文不把它单列为第四个不变量,因为在
DAG
语境下它表现为``正确作者的区块终被传播且可用'',已被数据可用性不变量与广播活性共同涵盖,并不构成去认证所转移的额外安全性负担。因此本文不宣称三不变量是
RBC
性质的``完备分解'',而只主张:它们覆盖了去认证分析所需的全部安全性/活性维度。

\subsection{4.2
不变量一:非等价性(non-equivocation)}\label{ux4e0dux53d8ux91cfux4e00ux975eux7b49ux4ef7ux6027non-equivocation}

非等价性要求:对任意验证者(作为区块作者)与任意轮次,至多存在一个被认证的区块。换言之,一个拜占庭作者无法让两个不同的区块在同一
(author, round) 位置上都获得认证。RBC
的一致性性质直接给出这一保证------若两个正确进程分别交付了来自同一发送者的消息,则两条消息相同
{[}5{]}, {[}9{]}。

非等价性对 DAG 共识至关重要。DAG
的排序规则建立在``所有正确验证者看到一致的
DAG''这一前提上;若某作者能就同一轮提供两个不同区块,不同验证者就可能基于不同的局部
DAG 推出不同的全序,安全性随之崩溃。可以说,非等价性是把``局部
DAG''提升为``全局一致对象''的黏合剂。

\subsection{4.3
不变量二:数据可用性(availability)}\label{ux4e0dux53d8ux91cfux4e8cux6570ux636eux53efux7528ux6027availability}

数据可用性要求:一个被引用的区块,其内容确实能够被需要它的验证者获取。RBC
在交付一个区块时,保证至少 f+1
个正确进程已经持有该区块的内容,因而任何正确进程日后总能向它们索取并取得该内容
{[}5{]}, {[}6{]}。

可用性与前一个不变量有一个本质区别。非等价性是一个关于``图结构是否一致''的逻辑性质------它可以由签名、轮次规则、引用关系来约束。可用性则是一个关于``某份数据在物理上是否真的被足够多的诚实节点持有''的性质------它不能仅靠规则``推理''出来。无论排序规则多么精巧,它都无法凭空让一份从未被传播的数据变得可获取。可用性最终要求一个物理事实:数据确已被持有,或至少在被提交前可被恢复。Cachin
与 Tessaro
的异步可验证信息分散(AVID)正是为这一物理事实提供保障的专门机制:它用纠删码把数据切片分散存储,使数据可在部分节点缺席时被重构
{[}6{]}。这一区别是本文后续分析的关键。

\subsection{4.4 不变量三:因果历史完整性(causal-history
integrity)}\label{ux4e0dux53d8ux91cfux4e09ux56e0ux679cux5386ux53f2ux5b8cux6574ux6027causal-history-integrity}

因果历史完整性要求:当一个验证者看到某个认证区块时,该区块所引用的全部祖先区块在结构上都是良构的------即都是已认证的、且引用关系本身正确无环的区块。须特别说明:它是一个\textbf{纯结构性质},只关乎``祖先在 DAG 结构上是否良构'',而\textbf{不包含}``祖先内容在物理上是否可被获取''------后者属于 §4.3 的数据可用性不变量。剥离之后,因果历史完整性是一个可仅凭本地结构检查验证的逻辑型性质。在 certified DAG 中,该不变量由认证机制连同引用规则共同维持 {[}31{]}, {[}32{]}。

\subsection{4.5
三个不变量的异质性}\label{ux4e09ux4e2aux4e0dux53d8ux91cfux7684ux5f02ux8d28ux6027}

第 4.2--4.4
节揭示了一个对全文至关重要的事实:认证所维持的三个不变量并不同质。非等价性与因果历史完整性本质上是结构一致性问题,可由密码学签名、轮次规则、引用关系、提交过滤等逻辑手段来约束;本文称之为逻辑型不变量。数据可用性则是一个物理传播问题,无法仅由规则导出;本文称之为物理型不变量。

这一逻辑型/物理型的二分可以用一个判准来刻画:一个不变量是逻辑型的,当且仅当它可以仅凭对协议消息的本地检查(签名验证、引用关系核对)被验证;它是物理型的,当且仅当它的验证最终依赖于``某份数据是否真实存在于某处''这一不可由本地检查推断的事实。按此判准,非等价性与因果历史完整性是逻辑型的,数据可用性是物理型的。这一异质性意味着:当认证被移除时,三个不变量``重新安家''的难易程度不会相同。本文将在下一章证明,正是这一不对称,决定了隐性代价的分布。

这一二分可被一项受控部署实证支持。在真实 Sui/Mysticeti n=7 部署上,沿区块负载轴(0 / 1 KB / 10 KB)各取 3 次重复(共 9 次运行),逻辑型一侧如二分所预言保持平直:eqChk/commit 与 causDec/commit 对负载的回归 $R^2 \leq 0.04$------逻辑型不变量的代价与负载正交。物理型一侧在可测区间内由 RPC 往返次数主导而非字节量主导(详见 §9 的范围限定)。

\begin{figure}
\centering
\pandocbounded{\includegraphics[keepaspectratio,alt={图 4-1 逻辑型不变量的每提交开销随区块负载的变化(eqChk/commit、causDec/commit;0 / 1 KB / 10 KB)。两条线在可测负载区间内平直(回归 R²≤0.04),印证逻辑型不变量的代价与数据体积正交。注:本图为单一 uncertified 协议(Sui/Mysticeti n=7)关键路径故障下的部署实测;负载上限受协议块大小限制(\textgreater10 KB 触发 TooManyTransactionBytes)。}]{figures/figure_4_1_payload_heterogeneity.png}}
\caption{\textbf{图 4-1}
逻辑型不变量的每提交开销随区块负载的变化(eqChk/commit、causDec/commit;0
/ 1 KB / 10 KB)。两条线在可测负载区间内平直(回归
R²≤0.04),印证逻辑型不变量的代价与数据体积正交。注:本图为单一 uncertified
协议(Sui/Mysticeti
n=7)关键路径故障下的部署实测;负载上限受协议块大小限制(\textgreater10 KB
触发 TooManyTransactionBytes)。}
\end{figure}

\section{第 5 章
SubRQ2:去认证后的保证迁移与隐性代价}\label{ux7b2c-5-ux7ae0-subrq2ux53bbux8ba4ux8bc1ux540eux7684ux4fddux8bc1ux8fc1ux79fbux4e0eux9690ux6027ux4ee3ux4ef7}

本章运用第 3 章的三分判据与 C1-C
双判据,逐一追踪三个不变量在去认证后的命运,并据此盘点隐性代价。

\subsection{5.1
工作假设与检验程序}\label{ux5de5ux4f5cux5047ux8bbeux4e0eux68c0ux9a8cux7a0bux5e8f}

基于第 4.5 节的异质性分析,本文提出一个待检验的工作假设
H:非等价性与因果历史完整性大概率落入 ① 平移,数据可用性大概率落入 ②
削弱+补偿。必须强调,H 是一个假设而非结论------本章的任务是用 C1-C
双判据去验证或推翻它,而不是预设它成立。这一``假设---检验''的表述方式,正是为了避免第
3.4 节所警惕的同义反复。检验程序如下:对每个不变量,先在 uncertified
协议中定位承接它的机制(若存在),再用判据(a)核验该机制是否产生可归因成本,用判据(b)核验协议正确性证明是否依赖它,最后据保证强度是否降级判定其落入
① 还是 ②。

\subsection{5.2 非等价性的命运:落入 ①
平移}\label{ux975eux7b49ux4ef7ux6027ux7684ux547dux8fd0ux843dux5165-ux2460-ux5e73ux79fb}

在 uncertified DAG 中,区块以尽力而为的方式广播,不再经过 RBC
认证。这意味着一个拜占庭作者确实可以在广播层发生等价(equivocation)------向不同验证者出示不同的区块。乍看之下,非等价性似乎被放弃了。

但更细致的考察表明并非如此。以 Mysticeti {[}43{]}
为例,等价被允许``发生'',却在提交规则中被``中和'':一个等价的区块无法聚集足够的支持,因而不会被提交;协议在
commit 逻辑中显式地处理等价情形。Adelie {[}45{]}
进一步指出,通过区块视图规则与关键区块规则的结合,一个拜占庭验证者在整个
epoch 内能够制造的等价次数被限制在 O(n)
量级。由此可见,非等价性这一不变量并未被消除------它只是从``广播层的事前认证''被平移到了``提交规则中的事后中和''。

为使这一判定可操作,须指认 Mysticeti 提交规则中承接非等价性的具体机制。Mysticeti 的提交规则要求一个区块获得至少 2f+1 个不同验证者的``投票''(通过对该区块的引用)方可被判定为已提交。当一个拜占庭作者对同一 (author, round) 槽位构造两个等价区块 X 与 X′ 时,每个正确验证者只会看到其中之一(因广播为尽力而为)。关键在于:由于正确验证者总数为 n−f ≥ 2f+1,而 X 与 X′ 的支持集之并至多覆盖全部 n 个验证者,在最坏情况下 X 获得 f+1 票、X′ 获得 f+1 票,但二者各自均无法达到 2f+1 的提交阈值(因 f+1 $<$ 2f+1 当 f $>$ 0)。因此,等价区块的两个分身都无法被提交,非等价性在提交层被``中和''。Mysticeti 的安全性论证(其 Lemma 4.2)明确依赖``每个 (author, round) 槽位至多一个区块被提交''这一引理,而该引理的成立正是上述支持计数机制的直接推论。

按 C1-C 判据:协议的安全性证明确实在某一步依赖``等价区块不会被提交''(判据 b
满足),而这一步对应着提交规则中可识别的逻辑复杂度(判据 a
满足)。上述 Mysticeti 的支持计数机制即判据(b)所要求的、``证明中明确依赖该不变量某版本''的可指认步骤。因此,非等价性落入
① 平移。

这一平移并非全无代价,但其代价是逻辑性的、可由规则界定的。等价的事后中和把原本由
RBC
统一承担的逻辑搬进了提交规则,使提交规则更复杂,却没有引入新的物理资源需求。这正是逻辑型不变量``易于重新安家''的体现。

\subsection{5.3 因果历史完整性的命运:落入 ①
平移}\label{ux56e0ux679cux5386ux53f2ux5b8cux6574ux6027ux7684ux547dux8fd0ux843dux5165-ux2460-ux5e73ux79fb}

因果历史完整性在 uncertified DAG
中由引用关系的验证与提交过滤共同维持。当一个区块被纳入提交时,协议要求其引用的祖先满足特定条件方可一并提交;不满足条件的引用会被规则过滤
{[}42{]},
{[}43{]}。这一机制以确定性规则替代了认证所提供的``祖先必为认证区块''的保证。同样按
C1-C 判据,该不变量落入 ① 平移:它由可识别的规则维持(判据 a、b
均满足),保证强度与认证版本基本相当。须明确:依第 4.4
节的界定,这里平移的只是因果历史完整性的\textbf{结构成分}(祖先良构、引用正确);祖先内容``是否可得''不在此不变量之内,而由下一节的数据可用性单独追踪。正因如此,因果历史完整性的低成本
① 平移并未``偷渡''任何物理可用性的代价------物理代价被完整地留给了
§5.4。

值得注意的是,因果历史完整性的平移与非等价性的平移共享同一个特征:二者都把保证从``广播层''搬到了``提交层'',而提交层本就是
DAG
共识需要逻辑处理的地方。逻辑型不变量之所以可低成本平移,根本原因在于它们的``新家''------提交规则------已经存在,平移只是往里面增添逻辑,而不需要新建一个机制。

非等价性与因果历史完整性同落 ①
平移这一判定,可由一项受控部署实证佐证。本人在真实 Sui n=7 与 n=4
部署上,沿独立丢包率(0 / 2\% / 5\% / 10\%)各取 5 次重复(共 40
次运行),直接测量两条逻辑型 enforcement
路径的活动量(eqChk、causDec)。结果显示:eqChk 与 causDec
随每次提交稳定发生(其计数与 commit 数同步),在 0--10\%
丢包全区间内从不被关断,cadence 亦维持基线(4.0--4.7
ckpt/s,无塌缩)------逻辑型不变量的维护是一笔在良好情形下即已恒定支付的固定每提交开销,符合判据(a)所定义的``可归因且非或然''的成本形态。这与下一节数据可用性的或然成本形成鲜明对照:同样这
40 次运行中恢复拉取(recFetch)恒为 0(见
§5.4),即逻辑型路径始终在线、而物理型补偿路径在纯丢包下从不触发。

\begin{figure}
\centering
\pandocbounded{\includegraphics[keepaspectratio,alt={图 5-1 两条逻辑型 enforcement 路径的每提交开销随丢包率的变化(eqChk/commit、causDec/commit;0--10\% 独立丢包,n=7)。两线平直、从不被关断,印证非等价性与因果历史完整性是 ① 平移后的固定每提交开销。注:本图为单一 uncertified 协议(Sui/Mysticeti)部署实测。}]{figures/figure_5_1_logical_paths_flat.png}}
\caption{\textbf{图 5-1} 两条逻辑型 enforcement
路径的每提交开销随丢包率的变化(eqChk/commit、causDec/commit;0--10\%
独立丢包,n=7)。两线平直、从不被关断,印证非等价性与因果历史完整性是 ①
平移后的固定每提交开销。注:本图为单一 uncertified
协议(Sui/Mysticeti)部署实测。}
\end{figure}

\subsection{5.4 数据可用性的命运:落入 ②
削弱+补偿}\label{ux6570ux636eux53efux7528ux6027ux7684ux547dux8fd0ux843dux5165-ux2461-ux524aux5f31ux8865ux507f}

数据可用性的情形与上述两者根本不同,这正是工作假设 H 的关键。

在 certified DAG
中,认证(RBC)在一个区块被引用之前就保证了其内容被足够多的诚实节点持有。在
uncertified DAG
中,这一事前保证消失了:一个区块可以在其内容尚未被充分传播时就被其他区块引用。如果该区块的作者是拜占庭的,它甚至可能让某个区块``看起来存在于
DAG 中''却始终不提供其内容。

uncertified
协议如何应对?它们并不去``恢复''事前的可用性保证------因为那恰恰需要
RBC,即恰恰是被移除的东西。它们采取的是一种弱化加补偿的策略:可用性被弱化为一个更弱的版本------``数据在被提交之前可被恢复'',而非``数据在被引用之前已被持有'';同时,协议引入补偿机制:在提交时检查数据是否可得,对不可得的区块执行跳过或重试。DispersedLedger
{[}47{]}
更明确地展示了这一思路:让节点先就区块的承诺达成一致,再各自按自身带宽异步地下载内容;Al-Bassam
等人关于欺诈与数据可用性证明的工作 {[}48{]}
也表明,可用性的保障在去除强广播原语后需要专门的、独立的机制(纠删码加抽样)来承接。

按 C1-C 判据,可用性显然没有落入 ①
平移:其保证强度被实实在在地降级了(从``引用前已持有''降为``提交前可恢复'')。它也没有落入
③ 消除:协议确实引入了补偿机制,这些机制在成本台账中可见(判据
a),并在活性证明中被依赖(判据 b)。因此,数据可用性落入 ②
削弱+补偿。工作假设 H 得到验证。

这一结论的深层原因,正是第 4.5
节的逻辑型/物理型二分。可用性是物理型的:它的``新家''不像提交规则那样现成存在------要真正保证数据被持有,必须有节点真实地存储与传输数据,这是一个无法由规则凭空创造的物理资源。uncertified
协议无法为可用性提供一个零成本的逻辑新家,只能弱化它,再用补偿机制兜底。

数据可用性落入 ②
削弱+补偿这一判定,在部署中表现得最为清晰,且其触发条件本身即是一项可分离的实证发现。本人在真实
Sui n=7 部署上做了两组互补的受控实验。其一(触发的阴性对照):在上节所述
0--10\% 独立丢包的 40 次运行中,补偿路径(按需恢复拉取 recFetch)恒为
0------随机丢包尽管推高重传,却不足以越过触发按需恢复的``悬崖'',因为缺失的区块总能经常规广播补达。其二(触发的阳性实证与剂量响应):当把故障由随机丢包改为网络分区(把少数验证者从
quorum 切离 30 / 60 / 90 s,3×3×2 共 18 次运行),recFetch
在每一次运行中均大于 0(最小 104,最大
1412)------断连(而非丢包)才是越过悬崖、迫使 ② 补偿机制 materialize
的条件。更进一步,在补偿路径中植入受控 Δ\_recover 旋钮后,愈合到 cadence
恢复的墙钟停顿 t\_resume 随 Δ\_recover 单调、近似线性上升(合并斜率 +20.0
s / 1000 ms,r=+0.71;gap=90 s 下 R²=0.99,gap=30 s 下 R²=0.92)------这正是
§5.5
所述``或有负债在故障情形下集中爆发''的直接部署剂量响应。而同样这些分区运行中,逻辑型路径(eqChk、causDec)从不被关断,印证两类不变量的成本形态截然不同:逻辑型恒定在线,物理型按需爆发。构念效度限定见
§7.6.2、§9(单一 uncertified 协议、Δ\_recover 为注入式旋钮)。

\begin{figure}
\centering
\pandocbounded{\includegraphics[keepaspectratio,alt={图 5-2 恢复拉取(recFetch)随独立丢包率的变化(0--10\%,n=7)。recFetch 在全部 40 次运行中恒为 0------随机丢包不足以越过触发按需恢复的悬崖。注:本图为单一 uncertified 协议部署实测,与图 5-3 的分区实验构成触发条件的阴/阳对照。}]{figures/figure_5_2_loss_no_pull.png}}
\caption{\textbf{图 5-2}
恢复拉取(recFetch)随独立丢包率的变化(0--10\%,n=7)。recFetch 在全部 40
次运行中恒为 0------随机丢包不足以越过触发按需恢复的悬崖。注:本图为单一
uncertified 协议部署实测,与图 5-3 的分区实验构成触发条件的阴/阳对照。}
\end{figure}

\begin{figure}
\centering
\pandocbounded{\includegraphics[keepaspectratio,alt={图 5-3 网络分区下或有负债的部署剂量响应(n=7)。左:每次分区运行 recFetch 均 \textgreater0(最小 104),与图 5-2 的丢包对照鲜明;右:恢复停顿 t\_resume 随补偿路径 Δ\_recover 旋钮近似线性上升(gap=90 s,R²=0.99)。注:本图为单一 uncertified 协议部署实测;Δ\_recover 为注入式旋钮。}]{figures/figure_5_3_partition_dose_response.png}}
\caption{\textbf{图 5-3}
网络分区下或有负债的部署剂量响应(n=7)。左:每次分区运行 recFetch 均
\textgreater0(最小 104),与图 5-2 的丢包对照鲜明;右:恢复停顿 t\_resume
随补偿路径 Δ\_recover 旋钮近似线性上升(gap=90 s,R²=0.99)。注:本图为单一
uncertified 协议部署实测;Δ\_recover 为注入式旋钮。}
\end{figure}

这一对比可通过同一区块 b 在两类协议中的提交路径进一步具象化。在 certified DAG 中,b 先经 RBC 获得证书,其内容已被 f+1 个正确节点持有;此后任何引用 b 的区块都可安全地假定 b 可用,提交时无需额外检查。在 uncertified DAG 中,b 被直接广播并可立即被引用;当某个锚点的提交牵连到 b 时,协议必须在此刻确认 b 的内容可得------若不可得,则触发恢复。去认证并未消除可用性的代价,只是改变了它被支付的时机与条件------而``按需支付''在数据确实缺失时会变得昂贵。

\subsection{5.5
隐性代价台账}\label{ux9690ux6027ux4ee3ux4ef7ux53f0ux8d26}

至此可以盘点隐性代价。去认证并非``删去一个开销模块'',而是把三个不变量的维护责任从认证这一单一机制,重新分配到了协议的多个其他组件。这一重新分配带来四类隐性代价:

第一,\textbf{提交规则复杂度上升}。非等价性的事后中和、因果历史的提交过滤,都把原本由
RBC 统一承担的逻辑塞进了 commit 规则,使其显著复杂化。Adelie {[}45{]}
为应对 uncertified DAG
的新攻击面而专门设计检测与防护机制,正是这一代价的直接体现。

第二,\textbf{最坏情况延迟与恢复成本}。可用性的补偿机制(跳过、重试、按需恢复)在数据确实缺失时会触发,触发时的延迟远高于良好情形。这一成本在良好情形下为零,却在故障情形下集中爆发------它是一笔``或有负债''。

第三,\textbf{故障恢复难度上升}。在 certified DAG
中,认证为故障恢复提供了稳固的锚点;uncertified DAG
失去这一锚点后,恢复逻辑必须自行处理``被引用但不可用''的区块,恢复路径的状态空间随之扩大。

第四,\textbf{安全性证明负担上升}。认证为协议正确性证明提供了简洁的引理;去认证后,证明必须显式地论证事后中和机制的正确性,论证负担随之转移并加重。这一点虽不直接影响运行时性能,却影响协议的可验证性与可信度。

\subsection{5.6 反例排除}\label{ux53cdux4f8bux6392ux9664}

本文的中心主张断言``去认证从不产生 ③
消除''。这是一个全称否定,严格意义上无法被穷举证明。本文采取两条策略来支撑它。其一(F1),对本文语料库
C 中的每一个 uncertified
协议,逐一核验其安全性与活性证明,确认证明的某一步必然依赖某个承接非等价性或可用性的机制;在
Cordial Miners {[}42{]}、Mysticeti {[}43{]}、Mahi-Mahi {[}44{]}、Adelie
{[}45{]}
中,这一核验均成立。以 Cordial Miners {[}42{]}为例,其 blocklace 结构虽弃用 RBC,但通过偏序数据结构的内在约束来 enforce 靇等价性:在 blocklace 中,一个验证者的每个新区块必须引用其自身此前的所有区块,形成一条不可分叉的局部链;若拜占庭作者试图等价,其两个分身将各自引用不同的前驱,从而在偏序中产生不可调和的冲突------这一冲突在排序阶段被检测并拒绝。因此,Cordial Miners 的非等价性 enforcement 从 RBC 的一致性性质平移到了 blocklace 的结构性质,符合 ① 平移。其二(F2),对一个假想的``删去认证且不加任何补偿''的协议,可显式构造一个等价攻击击穿其安全性(详见第
6 章)。两条策略都不能把结论提升为普适定理,因此本文明确将主张限定在语料库
C 内:在 C
中,每一次去认证都是结构性替换,而非消除。下一章将针对其中的非等价性这一维度,给出一个不依赖逐协议归纳的、范围明确界定的补充论证------它不把整体主张``提升为定理'',但为其最易被质疑的一环提供一个比归纳更硬的支点。

\section{第 6 章
非等价性的必要性论证}\label{ux7b2c-6-ux7ae0-ux975eux7b49ux4ef7ux6027ux7684ux5fc5ux8981ux6027ux8bbaux8bc1}

第 5 章的结论建立在对语料库 C
中若干真实协议的归纳之上。本章针对其中的非等价性维度,补充一个不依赖逐协议归纳的论证。须先界定本论证的范围与作用:它\textbf{只覆盖三个不变量中的非等价性这一个},而非等价性恰是第
5
章判定为``低成本平移''的那一个;因此本章并不把中心主张整体``提升为定理'',它做的是一件更有限也更扎实的事------为中心主张最容易被质疑的一环(``去认证之后,非等价性是否真的还\textbf{必须}存在'')提供一个比归纳更硬的支点。

\subsection{6.1
模型、引理与命题}\label{ux6a21ux578bux5f15ux7406ux4e0eux547dux9898}

本文遵循``在范围明确的模型里严格论证、再诚实讨论推广''的策略。论证在纯异步模型中进行:n≥3f+1
个验证者,其中至多 f
个为拜占庭,网络异步(消息延迟无上界但最终送达),敌手可调度消息顺序。协议以轮次化
DAG 为通信结构,每个验证者的提交决定是其本地 DAG
的确定性函数;提交一个区块/锚点需要后续轮次中达到法定人数(quorum,规模
n−f,在最优容错 n=3f+1 时即 2f+1)的引用支持------这一 quorum
设定涵盖语料库 C 中的全部协议。

本章论证的对象不是``所有可设想的 DAG-BFT
协议'',而是一类\textbf{安全性证明依赖槽位唯一性引理}的协议。所谓\textbf{槽位唯一性引理},指:``在任意正确验证者的本地
DAG 中,每个 (author, round) 槽位至多被一个区块占据。''
这一引理(或其等价形式)是 DAG-BFT
安全性证明的标准构件:正是它,才使排序规则可以把``槽位''当作全序中的唯一对象来处理。本文已逐一核验:语料库
C 中的全部协议(DAG-Rider、Narwhal/Tusk、Bullshark、Cordial
Miners、Mysticeti、Mahi-Mahi、Adelie)的安全性论证都使用了槽位唯一性引理,或与之等价的``每槽位至多一个被采纳区块''性质。

\begin{quote}
\textbf{命题 1}:设 P 为一个 DAG-BFT
协议,其总序安全性证明依赖槽位唯一性引理。则 P 必须包含某种 enforce
非等价性的机制;若不含,则该安全性证明不成立。
\end{quote}

\subsection{6.2
论证:槽位唯一性的破坏}\label{ux8bbaux8bc1ux69fdux4f4dux552fux4e00ux6027ux7684ux7834ux574f}

下文给出命题 1 的论证。须就严格度作一说明:本文以``论证 /
证明梗概''的标准呈现它------其各步在逻辑上是完整的,但对异步消息调度的完全形式化刻画(如
I/O 自动机级别的执行建模)超出本文范围。

采用反证法。假设 P 的安全性证明依赖槽位唯一性引理,但 P 不包含任何
enforce 非等价性的机制。

\textbf{第一步:等价区块可同时进入正确验证者的本地 DAG。} 由于 P
不阻止等价,一个拜占庭作者 b 可以构造两个不同区块 X≠X′,二者位于同一
(author, round) 槽位。b 把 X 发送给一部分正确验证者、把 X′
发送给另一部分。在异步网络中,敌手可自由调度消息送达顺序,使得在某个时刻确实存在两个正确验证者
v₁、v₂:v₁ 的本地 DAG 在槽位 (b,r) 上是 X,v₂ 的本地 DAG 在同一槽位上是
X′。这意味着槽位唯一性引理的结论为假------同一槽位在不同正确验证者的本地
DAG 中被不同区块占据。

\textbf{第二步:安全性证明因此失效。} P
的安全性证明依赖槽位唯一性引理,而该引理已被第一步证伪。需要具体指出它在证明链条中的哪一步失效:DAG-BFT
安全性证明的核心是 \textbf{quorum 交集论证}------任意两个规模为 n−f
的提交 quorum,其交集规模至少为 2(n−f)−n = n−2f;在 n≥3f+1 时 n−2f ≥
f+1,因而交集至少包含一个正确验证者;该正确验证者对每个槽位只持有唯一区块,于是它``钉住''了两个
quorum
在该槽位上的取值,迫使二者一致。这一步的有效性\textbf{前提性地依赖}槽位唯一性:它默认``任一正确验证者都全局性地为每个槽位钉住唯一区块''。一旦
b 等价,槽位 (b,r) 在不同正确验证者的本地 DAG
中已不再唯一,这一默认随之失效------交集中的正确验证者固然自身视图一致,但它对槽位
(b,r) 所持的区块只是 b 的两个分身(X 或 X′)之一,已无法迫使两个分别支持
v₁、v₂ 提交的 quorum 在该槽位上取值一致。quorum
交集不再蕴含取值一致,安全性证明的这一关键步骤断裂。

因此,在 P 不 enforce 非等价性的前提下,P
不存在一个依赖槽位唯一性的有效安全性证明------与``P
的安全性由该证明确立''矛盾。命题 1
得证:\textbf{安全性论证依赖槽位唯一性的 DAG-BFT 协议,必须 enforce
非等价性。}

\textbf{论证的范围:一个诚实的边界。} 命题 1 没有证明``一切可设想的安全
DAG-BFT 协议都必须 enforce
非等价性''这一与模型无关的强命题------原则上,可能存在一个用完全不同的证明技术(不经由槽位唯一性)建立安全性的协议,本文不排除这种可能。命题
1 证明的是一个\textbf{条件性}结论:只要协议沿用 DAG-BFT
标准的、以槽位唯一性为基石的安全性论证,非等价性就必须被
enforce。这一范围恰好是本文所需:已核验语料库 C
中的全部协议都属于这一类。这一论证与 FLP {[}2{]} 及 Bracha {[}5{]}
所揭示的异步环境根本约束一脉相承------它本质上是``异步网络中无法靠时间区分一致视图与分裂视图''这一困难,在
DAG 槽位结构下的具体化。

\textbf{论证的形式化边界。} 须明确命题 1 当前的严格度。本文以``论证/证明梗概''的标准呈现它------其各步在逻辑上是完整的,但对异步消息调度的完全形式化刻画(如 I/O 自动机级别的执行建模)超出本文范围。``槽位唯一性引理''在不同协议中的精确形式可能有所不同:在 DAG-Rider 中它源于 RBC 的一致性性质,在 Mysticeti 中它源于支持计数机制(详见 §5.2),在 Cordial Miners 中它源于 blocklace 的偏序约束。命题 1 的条件前提(``安全性证明依赖槽位唯一性'')对这些不同形式均成立,但本文不宣称已给出一个统一的、跨所有可能形式的形式化。此外,§6.5 的部分同步推广仅作论证性讨论,未给出逐字严格的证明------这是一个明确的边界,留作未来工作。

值得对照的是,本文在此\textbf{主动放弃}了一个更直观但不严格的论证:即``把正确验证者分成两组、让两组各自独立提交冲突全序''。该论证在最优容错
n=3f+1 下并不成立------两个能各自独立达到提交 quorum
的不相交正确节点组,需要至少 2(n−f) 个节点的参与,而
2(n−f)\textgreater n(因
n\textgreater2f),故无法构造。本文的论证不依赖``两组独立提交'',而只依赖``等价破坏了安全性证明所依赖的槽位唯一性引理'',从而绕开了上述算术障碍。

\subsection{6.3
必要性的另一面:问责性视角的佐证}\label{ux5fc5ux8981ux6027ux7684ux53e6ux4e00ux9762ux95eeux8d23ux6027ux89c6ux89d2ux7684ux4f50ux8bc1}

命题 1 从安全性角度论证了非等价性 enforcement
的不可或缺。问责性研究从另一个角度提供了佐证。Polygraph {[}51{]}
是首个可问责的拜占庭共识算法:当分歧发生时,它能够检测出制造分歧的恶意节点。其关键发现是,仅靠检测等价(equivocation)本身并不足以问责------问责需要至少
Ω(n) 轮的额外日志。Civit 等人随后的 ABC 框架 {[}52{]}
则表明,把任意拜占庭共识协议转化为可问责协议,需要额外两轮全体通信与 O(n²)
比特的开销。这些结果从代价侧印证了本文的论点:与等价相关的保证,无论是``事前阻止''还是``事后问责'',都必然伴随可观的、不可消去的成本。换言之,非等价性不仅在逻辑上必要(命题
1),其 enforcement 在代价上也不可能免费。

\subsection{6.4 一个推论:enforcement
时机的自由,消除的不自由}\label{ux4e00ux4e2aux63a8ux8bbaenforcement-ux65f6ux673aux7684ux81eaux7531ux6d88ux9664ux7684ux4e0dux81eaux7531}

命题 1 不限定 enforcement 以何种方式进行。它既可以是 certified DAG
中``引用前的 RBC 认证'',也可以是 uncertified DAG
中``提交前的等价中和''。协议在 enforcement
的时机与机制上是自由的------但在``是否
enforce''上没有自由。这正好解释了第 5.2
节的经验观察:非等价性只能被平移,不能被消除。去认证改变的是 enforcement
的位置,而非其存在性。

须重申命题 1
在全文论证中的确切位置:它只锁定\textbf{非等价性}这一维,而非等价性恰是三个不变量中第
5 章判定为``低成本 ① 平移''的那一个。因此命题 1 \textbf{不蕴含}``去认证
= 结构性替换''这一整体主张------后者还需第 5 章对数据可用性(②
削弱+补偿)与因果历史完整性(①
平移)的独立分析共同支撑;真正承载隐性代价的可用性维度,其论证依据的是第
5.4 节的归纳与第 7 章的成本结构,而非命题 1。命题 1
的作用因此是精确而有限的:它堵死了``去认证可以干脆消除非等价性''这条最诱人的退路,使中心主张的三分之一获得了归纳之外的支撑;另外三分之二仍依赖第
5 章的归纳性分析。

\subsection{6.5
向部分同步的推广:一个诚实的边界}\label{ux5411ux90e8ux5206ux540cux6b65ux7684ux63a8ux5e7fux4e00ux4e2aux8bdaux5b9eux7684ux8fb9ux754c}

命题 1 在纯异步模型中陈述。真实部署的 DAG-BFT(如 Bullshark
{[}33{]}、Mysticeti {[}43{]})多运行于部分同步模型。部分同步并不削弱命题
1 的实质:在 GST 之前,网络实际上是异步的,敌手仍可如 §6.2
第一步那样,把等价区块送入不同正确验证者的本地 DAG,从而在 GST
前的窗口内破坏槽位唯一性。而安全性是一个必须\textbf{始终}维持、不能``等到
GST 之后''才成立的性质------因此,协议若要在 GST
前的异步窗口内保持安全,仍必须 enforce 非等价性。

需要诚实标注的是:GST
之后的同步期为协议提供了额外的时间结构,一个完整严格的部分同步版本论证应针对``GST
前异步窗口''重新表述其消息调度假设。本文不宣称已给出部分同步下逐字严格的证明,这是一个明确的边界。但由于命题
1
已被界定为一个\textbf{条件性}命题(以``安全性证明依赖槽位唯一性''为前提),它本身并不依赖纯异步假设的全部强度------只要存在一个网络行为如异步的窗口、且协议在该窗口内需维持安全,论证即可适用。语料库
C 中的协议均满足这一条件。

\section{第 7 章
SubRQ3:延迟优势的条件化结论}\label{ux7b2c-7-ux7ae0-subrq3ux5ef6ux8fdfux4f18ux52bfux7684ux6761ux4ef6ux5316ux7ed3ux8bba}

命题 1 表明(在第 6 章界定的协议类内)非等价性必须被 enforce,第 5
章表明可用性必须被弱化补偿。本章把这些代价计入账本,回答
SubRQ3:uncertified 的延迟优势在何种条件下成立。

\subsection{7.1
成本台账方法}\label{ux6210ux672cux53f0ux8d26ux65b9ux6cd5}

本文不依赖任何单一来源报告的 benchmark 数字------按 SubRQ0
的警觉,这些数字本身可能是混淆的。本文采用成本台账方法:把协议运行中的通信步与每区块开销,逐项归因到``维持某个不变量的某个机制''。这样得到的不是一个标量延迟,而是一个按不变量与情形分解的成本结构。

须强调一项记账原则:公平的成本台账必须让 certified 与 uncertified
双方的成本都随故障与不利网络而退化------certified DAG 在故障下同样有 RBC
重传、视图切换等开销,并非一个固定不变的基线。因此本章关注的不是某一方的绝对最坏情况,而是两类设计的成本随条件恶化的\textbf{相对斜率}:uncertified
一侧的或有负债,是否比 certified
一侧增长得更快。下文的``翻转''判断,指的正是相对斜率的交叉,而非把一方的最坏情况与另一方的良好情况直接对比。

\subsection{7.2
一个简化的成本模型}\label{ux4e00ux4e2aux7b80ux5316ux7684ux6210ux672cux6a21ux578b}

设 uncertified 版本相对认证版本在每轮节省的延迟为 Δ\_save(主要来自省去
RBC 的额外通信步),补偿机制每次触发的额外延迟为
Δ\_recover,补偿机制被触发的频率为 ρ(0≤ρ≤1)。则 uncertified 相对
certified 的净延迟优势可近似表示为

净优势 ≈ Δ\_save − ρ · Δ\_recover。

这一模型有意保持简化,其目的不是给出精确数值,而是揭示三个变量之间的定性关系。两个观察立即可得。其一,当
ρ→0(良好情形)时,净优势 ≈ Δ\_save\textgreater0,uncertified 占优;这与
Cordial Miners 把良好情形延迟由 4 步降至 3 步、Mysticeti 逼近 3
个消息轮次下界的报告一致 {[}42{]}, {[}43{]}。其二,在把 ρ 与 Δ\_recover
视为相互独立这一\textbf{显式建模假设}下,可解出一个临界频率 ρ* = Δ\_save
/ Δ\_recover,当 ρ\textgreater ρ*
时净优势转负------延迟优势被补偿成本反向吞噬。须强调这一独立性是一个被有意识采纳的简化:第
7.4 节将指出真实系统中 ρ 与 Δ\_recover 可能相互耦合,届时 ρ*
不再是一个良定义的标量阈值,而应被理解为``存在一条净优势归零的边界''这一定性事实。本文在用到
ρ*
时,凡涉及精确数值即默认处于该独立性假设内,凡涉及定性结论(存在翻转)则不依赖它。§7.6.2
的部署标定为这一独立性简化提供了一个旁证:Δ\_recover(每事件恢复时延)与恢复事件频率作为两个成本轴在实测中可分离------§7.6.2 的旋钮标定显示,当旋钮增大时单次观测的恢复事件计数反而下降,而总 cadence 代价仍上升,即代价由``每事件时延''而非``事件次数''承载,二轴可分离。独立性假设的适用边界可据此定量界定:在 $\Delta$\_recover $<$ 1000 ms 且 $\rho$ $<$ 0.04 次/commit 的范围内(即 §7.6.2--7.6.3 的实测覆盖区间),二轴的可分离性获实测支持;超出此范围(如 $\Delta$\_recover $>$ 2000 ms 的极端恢复成本),§7.6.3 的 amortized 观测已显示净优势对 $\Delta$\_recover 呈非单调($R^2 \approx 0.02$),独立性假设不再适用,此时 $\rho$* 应被理解为``存在一条净优势归零的边界''这一定性事实,而非一个良定义的标量阈值。§7.6.3
进一步以含 certified 对照臂的部署实测,把这一翻转频率 ρ* 首次定量定位在约
0.013--0.038 次/commit(单引擎可用性证书操作化下),并诚实地标出 amortized
观测对线性闭式的偏离。下文逐一分析推高 ρ(及 Δ\_recover)的三类条件。

为核对闭式模型的实现正确性,本文用一个 Monte Carlo 仿真(2000 轮、500 条 trial)展开为逐轮随机过程;仿真均值在全部参数点上收敛到解析值(相对误差 $<$ 0.5\%),验证闭式公式被正确实现。须明确这一核对属于实现层面的交叉核对,\textbf{不}构成模型贴合真实协议的证据------后者由 §7.6.3 的部署标定提供。图 1 给出该模型在 $\Delta$\_save = 2 时的净优势曲线。

\begin{figure}
\centering
\pandocbounded{\includegraphics[keepaspectratio,alt={图 1 uncertified 相对 certified 的净延迟优势随恢复触发频率 ρ 的变化。Δ\_save = 2;三条曲线对应不同的单次恢复成本 Δ\_recover;圆点标出净优势归零的临界频率 ρ*。注:本图为本节解析成本模型的数值图示,非真实协议实测。}]{figures/figure_1_net_advantage.png}}
\caption{\textbf{图 1} uncertified 相对 certified
的净延迟优势随恢复触发频率 ρ 的变化。Δ\_save =
2;三条曲线对应不同的单次恢复成本 Δ\_recover;圆点标出净优势归零的临界频率
ρ*。注:本图为本节解析成本模型的数值图示,非真实协议实测。}
\end{figure}

\subsection{7.3
翻转条件一:恢复频率}\label{ux7ffbux8f6cux6761ux4ef6ux4e00ux6062ux590dux9891ux7387}

第 5.6
节指出,可用性补偿机制的成本是一笔或有负债:良好情形下为零,故障情形下集中爆发。当系统中拜占庭节点积极地利用
uncertified DAG
的新攻击面------例如故意引用不提供内容的区块------补偿机制的触发频率 ρ
将显著上升。Adelie {[}45{]} 为应对正是这类新攻击面而引入额外机制,Shoal++
{[}38{]} 与 Autobahn {[}49{]}
对故障与不利网络条件下延迟的专门关注,都印证了故障情形并非可忽略的边缘情形,而是协议设计必须正视的常规情形。尽管本文不依赖任何单一来源的精确数字,这些已发表的故障实验为
ρ 的量级提供了定性锚点:Shoal++ {[}38{]} 与 Autobahn {[}49{]}
均专门测量了故障与网络抖动条件下的延迟,并报告了相对良好情形的显著劣化------这表明在真实部署中
ρ 远非可忽略的零,``故障情形''在成本台账中应作为常规列项,而非边角修正。当
ρ 越过临界值 ρ*,uncertified 的延迟优势消失。

翻转条件一可被一项受控部署实证直接标定。本人在真实 Sui n=7
部署上固定故障数(N=2 次关键路径轮转故障),通过调节故障间距(settle ∈ \{20,
40, 80, 160\} s)扫出一个真实的四档恢复频率梯度 ρ ∈ \{0.0077, 0.0057,
0.0037, 0.0022\} 次/commit(每档 3 次重复,共 12 次运行),直接测量共识
cadence 随 ρ 的退化。结果给出一条干净的负斜率直线:cadence = 3.78 −
265·ρ(n=4 个数据点,Pearson r=−0.95,$R^2$=0.90,斜率 95\% CI $[$−429, −101$]$,p=0.025)------恢复频率每升高,吞吐即线性下行,这正是
§7.2 成本模型中 ρ·Δ\_recover 项在部署上的体现:ρ
越大,或有负债被触发得越频繁,净优势被吞噬得越多。配套地,原始 ρ
扫描(故障数
N=0\ldots8)给出每故障恢复成本的近似定值``量子'':recFetch≈143--164 /
故障、t\_resume≈37--41 s / 故障(两次独立扫描互相吻合),而无故障基线
cadence(4.32 ckpt/s)在任何故障注入下即跌至约 1.6 ckpt/s------这为 §7.2
中``每次触发的定额恢复成本''提供了部署量纲。须明确构念效度限定(同
§7.6.2):实证在单一 uncertified 协议(Sui/Mysticeti n=7)上展开、无
certified 对照臂,Δ\_recover 为注入式旋钮;故本段标定的是 uncertified 一侧
cadence 随 ρ 的退化斜率,而非 certified/uncertified 的净优势翻转点(后者见
§7.6.3)。

\begin{figure}
\centering
\pandocbounded{\includegraphics[keepaspectratio,alt={图 7-6 共识 cadence 随恢复频率 ρ 的退化(n=7,N=2 固定,经故障间距扫出四档 ρ)。拟合 cadence=3.78−265·ρ(R²=0.90),恢复频率升高即线性压低吞吐。注:本图为单一 uncertified 协议部署实测;Δ\_recover 为注入式旋钮,无 certified 对照臂(净优势翻转 ρ* 见 §7.6.3)。}]{figures/figure_7_6_rho_degradation.png}}
\caption{\textbf{图 7-6} 共识 cadence 随恢复频率 ρ 的退化(n=7,N=2
固定,经故障间距扫出四档 ρ)。拟合
cadence=3.78−265·ρ(R²=0.90),恢复频率升高即线性压低吞吐。注:本图为单一
uncertified 协议部署实测;Δ\_recover 为注入式旋钮,无 certified
对照臂(净优势翻转 ρ* 见 §7.6.3)。}
\end{figure}

\subsection{7.4
翻转条件二:部分同步退化与拥塞反馈}\label{ux7ffbux8f6cux6761ux4ef6ux4e8cux90e8ux5206ux540cux6b65ux9000ux5316ux4e0eux62e5ux585eux53cdux9988}

uncertified
协议的低延迟在很大程度上依赖良好的网络同步期。当系统进入部分同步退化(频繁的
GST 抖动)时,等价中和与可用性恢复都更频繁地被触发,Δ\_recover 与 ρ
同时上升。Autobahn {[}49{]} 明确指出,传统 BFT
协议在网络抖动后会产生``宿醉''(hangover)------积压的请求导致延迟在同步恢复后仍长期劣化。uncertified
DAG
的补偿机制在拥塞反馈下同样可能陷入这种正反馈式的延迟劣化:恢复请求本身占用带宽,带宽紧张又推高数据缺失率,从而进一步推高
ρ。须区分两类效应:本节所述 ρ 与 Δ\_recover
在退化下的``同向上升''是二者受共同驱动(网络退化)的协变,而非
Δ\_recover→代价映射本身的非线性;§7.6.2
的受控旋钮标定显示,在固定故障强度下 Δ\_recover
轴的传导近似线性、非线性集中于 ρ
轴,且或有负债仅在恢复处于共识关键路径时才转化为可观测代价。

翻转条件二在部署中给出一个诚实的阴性结果,而这一阴性本身具有信息量。本人在真实
Sui n=7 部署上注入周期性 sub-critical GST 抖动(每 30 / 60 / 90 s
把少数验证者 \{v1, v2\} 短暂分区 5 s,每档 3 次重复,共 9 次运行),测量
cadence 是否随抖动频率退化。结果:cadence
在三档周期下均维持近基线(4.14--4.38 ckpt/s),cadence 对抖动周期的回归
R²=0.09(近乎无关)------sub-critical
抖动被协议吸收,未引发级联停顿。这一阴性结果精化(而非削弱)§7.4
的论断:并非任何部分同步扰动都会触发或有负债;只有当扰动严重到把恢复推上共识关键路径(如
§7.6.2 的轮转故障所构造)时,Δ\_recover 与 ρ
的同向上升才转化为可观测的吞吐代价。换言之,翻转条件二的成立有一个强度门槛,与
§7.6.1 沿故障强度轴的发现一脉相承。构念效度限定同 §7.6.2(单一
uncertified 协议、注入式旋钮;此为诚实阴性结果,sub-critical
设计下信号本就弱)。

\begin{figure}
\centering
\pandocbounded{\includegraphics[keepaspectratio,alt={图 7-7 cadence 随 sub-critical GST 抖动周期的变化(n=7,周期 30 / 60 / 90 s,少数节点 5 s 分区)。cadence 维持近基线(回归 R²=0.09)------sub-critical 抖动被吸收,印证或有负债的触发存在强度门槛。注:本图为单一 uncertified 协议部署实测;此为诚实阴性结果。}]{figures/figure_7_7_gst_jitter.png}}
\caption{\textbf{图 7-7} cadence 随 sub-critical GST
抖动周期的变化(n=7,周期 30 / 60 / 90 s,少数节点 5 s 分区)。cadence
维持近基线(回归 R²=0.09)------sub-critical
抖动被吸收,印证或有负债的触发存在强度门槛。注:本图为单一 uncertified
协议部署实测;此为诚实阴性结果。}
\end{figure}

\subsection{7.5
翻转条件三:网络非对称与负载偏斜}\label{ux7ffbux8f6cux6761ux4ef6ux4e09ux7f51ux7edcux975eux5bf9ux79f0ux4e0eux8d1fux8f7dux504fux659c}

DispersedLedger {[}47{]}
表明,验证者带宽的差异会显著影响共识表现。在网络非对称或负载偏斜的条件下,uncertified
DAG
的``按需恢复''策略会让低带宽节点成为新的瓶颈------它们既要参与排序,又要在数据缺失时承担恢复开销,从而推高其感知到的
Δ\_recover。认证版本由于在传播阶段就均摊了数据分发,反而在这类条件下更稳健。这揭示了一个反直觉的事实:认证的``事前均摊''在不利网络下是一种鲁棒性资产,而非单纯的开销。这一论断需要一个诚实的限定:RBC
自身的 O(n²) 消息复杂度在带宽非对称下也会压迫低带宽节点,certified
一侧并非没有自己的非对称惩罚。本文的主张因此是\textbf{相对的}------事前、均匀支付的可用性成本,其随网络非对称恶化的斜率,通常小于``按需恢复''这种在数据缺失时集中爆发的成本;而非声称
certified 在不利网络下绝对占优。这正是 §7.1
所述``比较相对斜率而非绝对最坏情况''记账原则的一次具体应用。

翻转条件三可被一项含 certified 对照臂的受控部署实证直接确证。本人在真实
Sui n=7 部署上,把全部验证者带宽限制到 1 Mbps(经容器内 tc tbf
施加)以制造非对称压力,并在关键路径轮转故障下对照 §7.6.3 所述可开关的
certificate-gate 两臂(cert ON / OFF),各取 3 次重复(共 9
次运行,另含对称带宽基线臂)。结果直接印证本节核心论断:在同等 1 Mbps
非对称下,cert ON 比 cert OFF 同时取得更高 cadence(1.96 vs 1.06
ckpt/s,+86\%)、更短恢复停顿(t\_resume 26.1 vs 44.4
s,−41\%)与更少恢复拉取(recFetch 97 vs
188,−48\%)------认证在传播阶段事前均摊数据分发,使落后节点恢复时其缺失祖先已被
quorum
持有,故在不利网络下恢复显著更廉。这是``认证的事前均摊在非对称网络下是鲁棒性资产''这一反直觉论断的首个部署确证。须保留本节既有的相对性限定:此处对照臂为
§7.6.3 的可用性证书 gate(2f+1 接受,非密码学签名聚合),单引擎、Δ\_recover
注入式;且 cert OFF 一侧有一次重复为停顿离群点(cadence
方差较大,±0.84),故可信的是优势方向与量级而非精确百分比。

\begin{figure}
\centering
\pandocbounded{\includegraphics[keepaspectratio,alt={图 7-8 1 Mbps 带宽非对称下 certificate-gate 两臂的对比(cert ON vs OFF,n=7,关键路径轮转故障)。三联图:cadence、t\_resume、recFetch;cert ON 在三项上均更优,印证事前均摊在不利网络下的鲁棒性。注:本图为真实 Sui 部署实测;cert-gate 为可用性证书(非密码学签名),Δ\_recover 为注入式旋钮;cert OFF 一次重复为离群点。}]{figures/figure_7_8_net_asymmetry.png}}
\caption{\textbf{图 7-8} 1 Mbps 带宽非对称下 certificate-gate
两臂的对比(cert ON vs
OFF,n=7,关键路径轮转故障)。三联图:cadence、t\_resume、recFetch;cert ON
在三项上均更优,印证事前均摊在不利网络下的鲁棒性。注:本图为真实 Sui
部署实测;cert-gate 为可用性证书(非密码学签名),Δ\_recover
为注入式旋钮;cert OFF 一次重复为离群点。}
\end{figure}

\subsection{7.6 条件化结论}\label{ux6761ux4ef6ux5316ux7ed3ux8bba}

综合 7.2--7.5,SubRQ3 的答案是:uncertified DAG
的延迟优势成立于一个明确界定的区域------网络接近同步、故障稀少、负载均衡(即
ρ\textless ρ*
的区域);在该区域之外,优势会被隐性代价抵消,在恢复频率足够高或网络足够非对称时甚至反转。所谓``结构性''的优势,实为一个条件性的优势。下表概括了这一条件化结论。

{\def\LTcaptype{} % do not increment counter
\begin{longtable}[]{@{}lll@{}}
\toprule\noalign{}
条件维度 & uncertified 占优 & 优势被抵消/反转 \\
\midrule\noalign{}
\endhead
\bottomrule\noalign{}
\endlastfoot
故障率 / 恢复频率 ρ & ρ 接近 0 & ρ 超过临界值 ρ* \\
网络同步性 & 长同步期、GST 稳定 & 部分同步退化、GST 频繁抖动 \\
带宽分布 & 验证者带宽对称 & 网络非对称、负载偏斜 \\
敌手行为 & 敌手不利用可用性攻击面 & 敌手主动制造数据缺失 \\
\end{longtable}
}

上表前三行的``优势被抵消/反转''条件现各有一项直接部署锚点:恢复频率 ρ 见
§7.3(cadence 随 ρ 线性退化,R²=0.90)、网络同步性见 §7.4(sub-critical
抖动被吸收的阴性结果及其强度门槛)、带宽分布见 §7.5(非对称下 certified
臂更鲁棒)。三者连同 §7.6.1--7.6.3
的成本模型标定,把本表从定性预测落为部署可观测(各锚点共享的单引擎/可用性证书/注入旋钮限定见
§9)。

图 2 把这一条件化结论可视化为一张参数空间地图:在 (ρ, Δ\_recover)
平面上,临界边界 ρ* 把平面切成两部分------其左上方为 uncertified
占优区,其右下方优势被抵消乃至反转。这张地图直观地显示,``uncertified
更快''并非平面上的普遍事实,而只是其中一个区域的局部事实。这就回答了贡献
B:uncertified 的延迟优势并非 fundamental,而是 conditional。

\begin{figure}
\centering
\pandocbounded{\includegraphics[keepaspectratio,alt={图 2 (ρ, Δ\_recover) 参数空间中的 uncertified 优势区域。颜色表示净延迟优势(\textgreater0 即 uncertified 占优);黑色实线为临界边界 ρ* = Δ\_save/Δ\_recover。注:Δ\_save = 2;本图为第 7 章解析成本模型的数值图示。}]{figures/figure_2_advantage_region.png}}
\caption{\textbf{图 2} (ρ, Δ\_recover) 参数空间中的 uncertified
优势区域。颜色表示净延迟优势(\textgreater0 即 uncertified
占优);黑色实线为临界边界 ρ* = Δ\_save/Δ\_recover。注:Δ\_save =
2;本图为第 7 章解析成本模型的数值图示。}
\end{figure}

需要强调,这一结论与``现有 benchmark 在说谎''不同。本文的论点更精确:现有
benchmark
多在良好情形附近(ρ≈0)采样,它们报告的数字在其采样区域内是真实的,但把一个高维性能函数在少数维度上的投影,误读为了整个函数的形状。这就回答了贡献
A。

\subsubsection{7.6.1
跨规模与跨强度部署实证}\label{ux8de8ux89c4ux6a21ux4e0eux8de8ux5f3aux5ea6ux90e8ux7f72ux5b9eux8bc1}

§7.6 的条件化结论可被一项受控部署实证定向支持。本人在真实 Sui
7-validator (n=7, f=2) 与 4-validator (n=4, f=1)
部署上,沿两条独立压力轴展开:故障数相对 BFT 容差 f(within f vs at/above
f)与单故障强度(2\% vs 5\% 丢包)。每个 cell 取 10--20 次受控重复(关键
cell 另以独立种子复测,合计 290+
次重复),故障注入采用最小成对设计:跨验证者(独立)或同验证者(叠加),测量共识
cadence 的 p95 真值。

{\def\LTcaptype{} % do not increment counter
\begin{longtable}[]{@{}
  >{\raggedright\arraybackslash}p{(\linewidth - 12\tabcolsep) * \real{0.1364}}
  >{\raggedright\arraybackslash}p{(\linewidth - 12\tabcolsep) * \real{0.1364}}
  >{\raggedright\arraybackslash}p{(\linewidth - 12\tabcolsep) * \real{0.1364}}
  >{\raggedright\arraybackslash}p{(\linewidth - 12\tabcolsep) * \real{0.1364}}
  >{\raggedright\arraybackslash}p{(\linewidth - 12\tabcolsep) * \real{0.1364}}
  >{\raggedright\arraybackslash}p{(\linewidth - 12\tabcolsep) * \real{0.1364}}
  >{\raggedleft\arraybackslash}p{(\linewidth - 12\tabcolsep) * \real{0.1818}}@{}}
\toprule\noalign{}
\begin{minipage}[b]{\linewidth}\raggedright
配置
\end{minipage} & \begin{minipage}[b]{\linewidth}\raggedright
n / f
\end{minipage} & \begin{minipage}[b]{\linewidth}\raggedright
故障数
\end{minipage} & \begin{minipage}[b]{\linewidth}\raggedright
强度
\end{minipage} & \begin{minipage}[b]{\linewidth}\raggedright
故障放置
\end{minipage} & \begin{minipage}[b]{\linewidth}\raggedright
可加性
\end{minipage} & \begin{minipage}[b]{\linewidth}\raggedleft
p95 比值
\end{minipage} \\
\midrule\noalign{}
\endhead
\bottomrule\noalign{}
\endlastfoot
Test 0 canonical & 7 / 2 & 2 (= f) & 2\% loss & 跨验证者 & additive &
1.03 \\
Test 1(修复后) & 7 / 2 & 2 (= f) & 2\% loss & 同验证者 & additive &
1.00 \\
n=4 ablation & 4 / 1 & 2 (\textgreater{} f) & 2\% loss & 跨验证者 &
additive & 0.97 \\
Test 2 & 7 / 2 & 2 (= f) & 5\% loss & 跨验证者 & super-additive &
1.08 \\
n=4 + 5\% loss & 4 / 1 & 2 (\textgreater{} f) & 5\% loss & 跨验证者 &
super-additive & 1.12 \\
\end{longtable}
}

\textbf{表 7-1} 跨规模/跨强度受控故障注入的可加性结果(p95 真
cadence;比值 = observed / expected\_additive,\textgreater1 即超可加)。

实证给出两项判断。其一,在中强度(2\%
丢包)下,跨验证者故障呈干净的可加性,且这一点与委员会规模无关------n=7(within
f,ratio 1.00--1.03)与 n=4(at/above f,ratio 0.97)一致:跨越 BFT 容差边界 f
本身不破坏可加性。其二,把单故障强度提高到 5\% 丢包,则在 n=7(ratio
1.08)与 n=4(ratio 1.12)上双双把 p95 推入超可加。换言之,§7.3
所述``或有负债在故障下集中爆发''在部署中表现为一笔由故障强度驱动且非线性(超可加)的尾部代价,其触发门槛是把验证者挤出
fast-quorum 参与的单故障强度,而非形式化的 BFT 安全边界 f。

须明确这一实证的范围:它在单一 uncertified 协议(Sui,Mysticeti
家族)上展开,没有 certified 对照臂;因此它探查的是 uncertified
一侧或有负债的形状,既非 §7.2 成本模型(Δ\_save、ρ*)的定量标定,亦非
certified/uncertified 的净优势比较。它为 §7.6
表格的``翻转''语言提供了具体的部署锚点------或有负债确为真实、非线性、强度驱动------但完整的跨协议受控标定仍属未来工作(见
§9、§10)。工具学与复现性限定见 §9``实证仪器局限'';完整 2×2 与逐 rep
数据见
\texttt{multivalidator/\allowbreak{}reports/\allowbreak{}stress\_\allowbreak{}test\_\allowbreak{}combined\_\allowbreak{}conclusion\_\allowbreak{}v3\_\allowbreak{}20260526.\allowbreak{}md}。

\begin{figure}
\centering
\pandocbounded{\includegraphics[keepaspectratio,alt={图 7-3 修正后的 2×2 稳定性网格:横轴为单故障强度(2\% / 5\% 丢包),纵轴为故障数相对 BFT 容差(within f / at/above f)。强度轴是唯一驱动超可加的轴;BFT 边界轴无独立效应。注:本图为部署实测数据的网格化呈现。}]{figures/figure_7_3_stability_grid.png}}
\caption{\textbf{图 7-3} 修正后的 2×2 稳定性网格:横轴为单故障强度(2\% /
5\% 丢包),纵轴为故障数相对 BFT 容差(within f / at/above
f)。强度轴是唯一驱动超可加的轴;BFT
边界轴无独立效应。注:本图为部署实测数据的网格化呈现。}
\end{figure}

\subsubsection{7.6.2
恢复侧或有负债的受控旋钮标定}\label{ux6062ux590dux4fa7ux6216ux6709ux8d1fux503aux7684ux53d7ux63a7ux65cbux94aeux6807ux5b9a}

§7.6.1 沿故障强度轴刻画了或有负债``是否爆发'';本节沿 §7.2 成本模型的
Δ\_recover
轴,直接标定``爆发时的代价随单次恢复成本如何缩放''。方法上,本人在同一 Sui
部署的共识核心(Mysticeti)按需恢复路径中植入一个受控延迟旋钮
Δ\_recover(每次 ancestor
拉取批次注入一段可配置时延),并设计一种轮转故障,把节点恢复置于共识关键路径上:先隔离一组验证者使其落后(网络在剩余
quorum 上继续推进而打开
gap),随即在愈合这组节点的同时隔离另一组,使当前在线集低于
quorum------此时全网必须等待落后节点完成按需拉取(受 Δ\_recover
门控)方能恢复出块。可观测量取``愈合到 cadence 恢复''的墙钟时延
t\_resume。

实证给出三项判断。其一(剂量响应):在 n=7 与 n=4 上,t\_resume 都随
Δ\_recover 单调、近似线性地上升(n=7:斜率 ≈ 14.5 s / 单位旋钮,Pearson
r=+0.85,n=12 次运行,斜率 p$<$0.01;n=4:≈ 19.0
s,r=+0.94,n=8 次运行,斜率 p$<$0.01),未见超线性------即在固定故障强度下,或有负债沿 Δ\_recover
轴是线性的。这与 §7.6.1``超可加由故障强度驱动''互补:非线性集中于
ρ(触发频率)轴,Δ\_recover
轴本身线性。一个机制旁证支持把二者作可分离轴建模:当旋钮增大时,单次观测的恢复事件计数反而下降,而总
cadence
代价仍上升------代价由``每事件时延''(Δ\_recover)而非``事件次数''(ρ)承载,二轴可分离。其二(可操作边界):同一旋钮在节点恢复不处于关键路径时几乎不可观测------单节点短暂离线的恢复被协议的推送式补块(peer
广播 backlog)绕过,t\_resume 对 Δ\_recover 的相关仅 r=+0.43
且非单调;只有当恢复被强制置于关键路径时,代价才干净显现。这把 §7.3
笼统的``故障情形''精化为一个可操作判据:或有负债仅在恢复落入共识关键路径时
materialize。其三(委员会结构依赖):n=4
的基线更低(更少节点、和解更快)但斜率更陡(更高的每单位 Δ\_recover
代价)------委员会越小、恢复冗余越少,Δ\_recover
的传导越直接,或有负债随委员会收缩而加剧。

须明确本节的范围与构念效度限定,其约束强于 §7.6.1。Δ\_recover
此处是一个注入式仿真旋钮(在恢复路径人为植入时延以模拟单次恢复成本),而非对某个
certified 协议真实恢复开销的测量;实证仍在单一 uncertified 协议上、无
certified 对照臂。因此本节标定的是或有负债沿 Δ\_recover
轴的形状与可操作触发边界,而非 §7.2 中 Δ\_save、ρ* 的数值,也不是
certified/uncertified 的净优势比较。复现 harness 与脚本见
\texttt{stage9\_\allowbreak{}cross\_\allowbreak{}protocol\_\allowbreak{}calibration/\allowbreak{}},逐 rep
数据(\texttt{phaseA\_\allowbreak{}criticalpath\_\allowbreak{}*.\allowbreak{}csv})见
\texttt{stage7\_\allowbreak{}deployment/\allowbreak{}multivalidator/\allowbreak{}data/\allowbreak{}runs/\allowbreak{}}。

\begin{figure}
\centering
\pandocbounded{\includegraphics[keepaspectratio,alt={图 7-4 恢复侧或有负债的部署标定。左:n=4 与 n=7 下 cadence 恢复时延 t\_resume 随 Δ\_recover 的剂量响应(均近似线性,跨规模复现);右:同一旋钮在关键路径上(干净剂量响应)与非关键路径上(单节点恢复,被推送式补块绕过、无干净响应)的对比。注:本图为关键路径轮转故障下的部署实测;Δ\_recover 为注入式仿真旋钮,非真实 certified 协议开销。}]{figures/figure_7_4_recover_calibration.png}}
\caption{\textbf{图 7-4} 恢复侧或有负债的部署标定。左:n=4 与 n=7 下
cadence 恢复时延 t\_resume 随 Δ\_recover
的剂量响应(均近似线性,跨规模复现);右:同一旋钮在关键路径上(干净剂量响应)与非关键路径上(单节点恢复,被推送式补块绕过、无干净响应)的对比。注:本图为关键路径轮转故障下的部署实测;Δ\_recover
为注入式仿真旋钮,非真实 certified 协议开销。}
\end{figure}

\subsubsection{7.6.3 含 certified
对照臂的净优势翻转标定}\label{ux542b-certified-ux5bf9ux7167ux81c2ux7684ux51c0ux4f18ux52bfux7ffbux8f6cux6807ux5b9a}

§7.6.1、§7.6.2 在单一 uncertified
协议上刻画了或有负债侧的形状;它们都明确``无 certified 对照臂、非 §7.2
的定量标定''。本节补上这条对照臂,把 §7.2
成本模型的三个量------Δ\_save、认证对恢复成本的缩减、以及净优势翻转频率
ρ*------首次提升为真实 n=7 部署的实测。

操作化(受控 minimal-pair):在同一 Sui/Mysticeti binary
上把认证做成一个可开关的 certificate-gate。cert ON
时,一个父轮祖先区块仅当被一个 2f+1 权益 quorum 独立接受(经 Sui 已有的
RoundProber
确认通道直接观测,而非靠后续轮的提案引用隐式推断)后,方可被新区块引用------这把
certified-DAG
中``区块须先成证书再被构建于其上''的独立确认轮显式加回;cert OFF 即原生
uncertified Mysticeti。两臂共用同一引擎、同一提交规则与 leader
调度,仅``认证强度''这一单变量不同,从而规避 §3.3 所警惕的跨引擎混淆。

实证给出四项判断。其一(Δ\_save 为正且可观):良好情形(无故障)下 cert ON
的共识 cadence 系统性低于 cert OFF(4.58→2.99 ckpt/s,约 0.116
s/commit)------这正是 §7.2 中 uncertified
省去的认证确认轮延迟,首次在部署上被测得为一个正的、可观的
Δ\_save。其二(认证缩减恢复侧或有负债):在与 §7.6.2
相同的关键路径轮转故障、相同 Δ\_recover 下,cert ON
恢复得更便宜------恢复拉取次数更少(123 vs 196),恢复停顿 t\_resume
更短(每故障省 excess≈3--9 s),且这一节省随单次恢复成本 Δ\_recover
增大而增大(线性拟合斜率为正,R²≈0.49)。机制:gate 只引用已被 quorum
持有的区块,故落后节点恢复时其缺失祖先已被 2f+1 节点持有,拉取更廉。这是
§7.5``认证在传播阶段事前均摊数据分发、因而在不利条件下恢复更鲁棒''一论断的部署确证。其三(净优势翻转
ρ* 的定量定位):把良好情形的 Δ\_save(0.116
s/commit)与每故障的恢复成本之差 excess(Δ\_recover)
结合,可解出净优势归零的临界故障率 ρ* ≈ 0.013--0.038
次/commit(约一次关键路径故障每 26--78 个 commit):故障频率高于此则
certified 占优,低于此则 uncertified 占优;ρ* 随 Δ\_recover
增大而下降(恢复越贵,certified 在越低的故障率即完成翻转)。这是 §7.2
条件化翻转------``存在一条净优势归零的边界
ρ*''------在真实部署上的首次定量定位。其四(诚实的模型偏离,呼应
§7.4):若改用单一定窗的 amortized cadence 直接测净优势,则净优势对
Δ\_recover 呈非单调(线性拟合
R²≈0.02)------一次性的恢复成本被测量窗稀释,且在极高 Δ\_recover 下
certificate-gate 会延长地排斥仍在恢复、尚未被 quorum
接受的落后节点,形成一个与``恢复更廉''相竞争的效应。这一偏离印证了 §7.4
所述 ρ 与 Δ\_recover 在极端条件下的耦合:闭式 net\_advantage 在 amortized
观测下并非严格线性,只有用分解观测(清洁的 Δ\_save 速率 +
清洁的每故障恢复成本差)方能干净标定 ρ*。

须明确本节的范围与构念效度限定,其约束承自 §7.6.2 并进一步具体化。此
certificate-gate 实现的是认证的\textbf{可用性证书}(被一个 2f+1 quorum
接受),\textbf{非}密码学意义的 2f+1 签名聚合;Δ\_save 因此是``等 quorum
接受再引用''这一独立确认轮的延迟,经真实 RoundProber
网络通道度量------它比一个纯注入式 sleep
更忠实(确实动用了一个独立的确认往返),但仍\textbf{不是}对真实
Bullshark/Narwhal RBC 签名开销的测量。实证仍在单一引擎(Mysticeti 加/去
gate)上展开,而非两套独立协议的 benchmark;Δ\_recover 仍为注入式旋钮;ρ*
为一个 band(2 次重复,Δ\_recover=500
处有一个噪声离群点),可信的是其量级而非精确值。因此本节标定的是
certified/uncertified 净优势的\textbf{成本结构}------Δ\_save
为正且可观、认证确实缩减恢复侧或有负债、且存在一个 ρ*
翻转------而非某一对真实协议的绝对延迟比较。复现 harness、Route A1
共识源码改动与结果报告(\texttt{PHASE1\_\allowbreak{}RESULTS.\allowbreak{}md}、\texttt{PHASE2\_\allowbreak{}RESULTS.\allowbreak{}md})见
\texttt{stage9\_\allowbreak{}cross\_\allowbreak{}protocol\_\allowbreak{}calibration/\allowbreak{}},逐 rep
数据(\texttt{phase2c\_\allowbreak{}excess\_\allowbreak{}sweep\_\allowbreak{}summary.\allowbreak{}csv} 等)见
\texttt{stage7\_\allowbreak{}deployment/\allowbreak{}multivalidator/\allowbreak{}data/\allowbreak{}runs/\allowbreak{}}。

\begin{figure}
\centering
\pandocbounded{\includegraphics[keepaspectratio,alt={图 7-5 含 certified 对照臂的净优势翻转标定。左:关键路径故障下两臂的恢复停顿 t\_resume 随 Δ\_recover 的剂量响应(上:cert OFF/uncertified;下:cert ON/certified;绿带=认证带来的恢复节省 excess);右:由 Δ\_save 与 excess(Δ\_recover) 解出的净优势翻转故障率 ρ*(Δ\_recover)------其上 certified 占优、其下 uncertified 占优。注:本图为真实 n=7 Sui 部署实测;cert-gate 为可用性证书(非密码学签名),Δ\_recover 为注入式旋钮。}]{figures/figure_7_5_certified_arm_crossover.png}}
\caption{\textbf{图 7-5} 含 certified
对照臂的净优势翻转标定。左:关键路径故障下两臂的恢复停顿 t\_resume 随
Δ\_recover 的剂量响应(上:cert OFF/uncertified;下:cert
ON/certified;绿带=认证带来的恢复节省 excess);右:由 Δ\_save 与
excess(Δ\_recover) 解出的净优势翻转故障率 ρ*(Δ\_recover)------其上
certified 占优、其下 uncertified 占优。注:本图为真实 n=7 Sui
部署实测;cert-gate 为可用性证书(非密码学签名),Δ\_recover 为注入式旋钮。}
\end{figure}

\subsection{7.7
面向设计者的判定程序}\label{ux9762ux5411ux8bbeux8ba1ux8005ux7684ux5224ux5b9aux7a0bux5e8f}

第 7.6
节的结论表对设计者仍偏抽象。为提升可操作性,本文据上述分析给出一个粗粒度的判定程序------它不计算精确数值,只把``看情况''拆解为三个可分别估计的子判断,从而把一个具体部署定位到
§7.6 表格的某一行:

\begin{enumerate}
\def\labelenumi{\arabic{enumi}.}
\tightlist
\item
  \textbf{估计敌手与故障强度。}
  目标部署中是否存在有动机、有能力主动制造数据缺失的敌手?预期拜占庭比例是接近
  f 上限,还是远低于?据此把恢复频率 ρ 粗分为``近 0 / 中等 / 高''。
\item
  \textbf{估计网络条件。} 同步期是否长而稳定、GST
  抖动是否罕见?验证者带宽是否大致对称?据此把单次恢复成本 Δ\_recover
  粗分为``低 / 高''。
\item
  \textbf{定位。} 若(ρ 近 0)且(网络近同步、带宽对称),落入 §7.6
  表格的``uncertified 占优''行;若(ρ 中等以上)或(网络非对称 / GST
  频繁抖动),落入``优势被抵消乃至反转''行。
\item
  \textbf{边界情形。} 若 ρ
  中等且网络条件中等,本文的定性模型给不出确定结论;此时应回退到针对该具体部署的受控实测------这正是第
  10 章所列未来工作。
\end{enumerate}

这一程序的价值不在于精确,而在于它把``延迟优势是否成立''从一个无法回答的笼统判断,转化为三个各自有据可依的工程估计;其输入(敌手模型、网络同步性、带宽分布)恰好都是部署方在选型阶段本就需要评估的量。

\section{第 8 章 主贡献:availability-enforcement
作为深层设计轴}\label{ux7b2c-8-ux7ae0-ux4e3bux8d21ux732eavailability-enforcement-ux4f5cux4e3aux6df1ux5c42ux8bbeux8ba1ux8f74}

前七章的分析指向一个统一的结论。本章把这一结论从全文的潜结构提升为显结构,给出本文的主贡献
C。

\subsection{8.1 从 SubRQ1--3
到一个深层轴}\label{ux4ece-subrq13-ux5230ux4e00ux4e2aux6df1ux5c42ux8f74}

回顾全文:第 4 章表明三个不变量异质,可用性是其中唯一的物理型不变量;第 5
章表明去认证后隐性代价集中于可用性;第 6 章表明非等价性必须被
enforce,故其 enforcement 时机才是变量;第 7
章表明延迟优势的条件,本质上由可用性补偿机制的成本结构(ρ、Δ\_recover)决定。这些结论汇聚成一个判断:\textbf{真正决定一个
DAG-BFT 协议行为的,不是``它是否认证'',而是``它如何 enforce
数据可用性''。} certified/uncertified
这个二分,只是这一深层变量的一个粗糙投影。

\subsection{8.2
深层轴的两个维度}\label{ux6df1ux5c42ux8f74ux7684ux4e24ux4e2aux7ef4ux5ea6}

本文把这一深层设计轴具体化为两个可操作的维度。

\textbf{维度一:availability-enforcement 时机。}
一个协议在何时强制保证一个区块的数据可用?本文区分三个位置:(i)引用前(pre-reference)------区块在能被其他区块引用之前,其可用性即被保证,certified
DAG
属此;(ii)提交前(pre-commit)------区块可被引用,但在被提交进入全序之前其可用性被检查;(iii)提交后(post-commit)------可用性仅在提交后经由按需恢复来保证。enforcement
时机越靠后,良好情形延迟越低(Δ\_save
越大),但故障情形的补偿成本越高(Δ\_recover 越大)。

\textbf{维度二:reconciliation trigger 条件。}
当数据缺失被检测到时,协调修复(reconciliation)在何种条件下、以何种粒度被触发?是每发现一个缺失区块即触发,还是批量触发?触发是否依赖超时?触发的代价由谁承担?这一维度决定了第
7 章模型中 ρ 与 Δ\_recover 的具体形态。

这两个维度共同构成的平面,才是描述 DAG-BFT
设计空间的恰当坐标系。把``是否认证''作为坐标,等于只观察了维度一的一个端点(引用前
vs 非引用前),并完全忽略了维度二。

\subsection{8.3
用深层轴重新定位现有协议}\label{ux7528ux6df1ux5c42ux8f74ux91cdux65b0ux5b9aux4f4dux73b0ux6709ux534fux8bae}

在这一坐标系中,现有协议不再是``两类'',而是散布的点。DAG-Rider
{[}31{]}、Narwhal/Tusk {[}32{]}、Bullshark {[}33{]}
位于维度一的``引用前''端点,reconciliation
几乎不被触发(因为认证已保证可用性)。Mysticeti {[}43{]} 与 Mahi-Mahi
{[}44{]} 位于``提交前/提交后''之间,reconciliation
由提交时的数据检查触发。DispersedLedger {[}47{]}
则是一个尤其有启发性的点:它把 enforcement
时机推后,却用纠删码编码的承诺把 reconciliation
的代价显著降低------这说明维度二上的精巧设计可以部分补偿维度一上的激进。也就是说,在本文的坐标系中,``低延迟''与``鲁棒性''不是非此即彼,而是可以通过在两个维度上分别调参来联合优化的------这是二分法完全看不到的设计自由度。图
3 把语料库中的协议定位在这一坐标系中:certified
协议聚集于横轴``引用前''端、纵轴低位,uncertified
协议散布于``提交前/提交后''区间,而 Adelie 与 DispersedLedger
等中间设计落在二分法无法安放、本文坐标系却能清晰定位的位置上。

\begin{figure}
\centering
\pandocbounded{\includegraphics[keepaspectratio,alt={图 3 深层设计轴上的协议定位。横轴为 availability-enforcement 时机(引用前 / 提交前 / 提交后),纵轴为单次 reconciliation 开销(定性)。注:本图为按第 8 章深层轴框架的定性定位,坐标不代表实测数据。}]{figures/figure_3_design_axis.png}}
\caption{\textbf{图 3} 深层设计轴上的协议定位。横轴为
availability-enforcement 时机(引用前 / 提交前 / 提交后),纵轴为单次
reconciliation 开销(定性)。定位依据为各协议论文中对 enforcement 时机与 reconciliation 机制的描述(§8.3 已详述);坐标为定性而非实测数据。}
\end{figure}

\subsection{8.4
深层轴的解释力}\label{ux6df1ux5c42ux8f74ux7684ux89e3ux91caux529b}

一个坐标系是否优于另一个,取决于它能否解释前者解释不了的现象。certified/uncertified
二分无法自然地安放 Adelie {[}45{]}:Adelie 构建在 uncertified DAG
之上,却又引入了一套接近``认证''功能的检测与防护机制------在二分法下它是一个尴尬的混合体,既不纯粹
certified 也不纯粹 uncertified。但在本文的深层轴上,Adelie
是一个清晰的点:它的 availability-enforcement 时机靠后(继承 Mysticeti 的
uncertified 结构),而其 reconciliation trigger
被设计得格外精细(借助区块视图规则限制等价次数,从而压低
ρ)。本文的坐标系因此能够解释二分法解释不了的中间设计,这正是它作为分析框架的解释力所在。

需要诚实说明并对贡献 C 的主张作一个收窄。贡献 C
的风险不在于``深层轴构造不出来''------它的两个维度已在前七章的分析中隐式贯穿;其真正风险在于被指为单纯的重命名。这一指控有其道理:维度一(availability-enforcement
时机)在很大程度上确实是 certified/uncertified
二分的连续化:其``引用前''端点就是 certified。因此本文明确收窄贡献 C
的主张:贡献 C
真正非平凡、不可被归为重命名的部分,是\textbf{维度二(reconciliation
trigger)的引入},以及由此带来的\textbf{``延迟与鲁棒性是二维平面上可联合优化的目标,而非一维取舍''这一设计自由度}。DispersedLedger
{[}47{]}
正是其证据:它在维度一上激进(时机靠后),却用维度二上的精巧设计(纠删码编码的承诺)把代价压低------这是任何只有维度一的坐标系都无法表达的。

进一步的证据来自语料库 C 中维度一相近但维度二各异的协议对。Adelie {[}45{]}与 Mysticeti {[}43{]}在维度一上位置相近(均继承 uncertified DAG 的提交前/提交后 enforcement),但维度二截然不同:Adelie 通过区块视图规则与关键区块规则主动限制拜占庭验证者在每个 epoch 内制造等价的次数(压低 $\rho$),而 Mysticeti 的维度二以逐缺失触发的按需恢复为主。这一差异在深层轴上清晰可见:二者在同一维度一位置上,因维度二设计不同而具有不同的鲁棒性特征------Adelie 的 $\rho$ 被显式压低,故其或有负债更难被触发。Certified/uncertified 二分法无法表达这一差异(二者同属 uncertified),深层轴却能。类似地,Cordial Miners {[}42{]}的 blocklace 结构在维度一上比 Mysticeti 更靠前(传播阶段即完成部分 enforcement),但其维度二的 reconciliation 机制与 Mysticeti 的按需恢复不同------blocklace 的偏序结构允许在传播阶段即检测冲突,而非推迟到提交阶段。这些案例共同表明:维度二(reconciliation trigger)是一个独立于维度一的设计变量,它在 certified/uncertified 二分法下被完全遮蔽。

概言之,贡献
C 的可辩护内核应被理解为:reconciliation trigger 是一个被
certified/uncertified 二分完全遮蔽的、独立的设计维度。其独立性判据为:若两个协议在维度一上位置相近、但在维度二上设计不同,且其成本结构($\rho$、$\Delta$\_recover)有显著差异,则维度二具有独立解释力。

一个坐标系的价值,还在于它能否指出尚未被占据的设计点。本文据此提出一个
\textbf{conjecture(尚待实现验证)}:``提交后 enforcement(维度一最靠后)+
批量触发的 reconciliation +
纠删码编码的可用性承诺(维度二)''这一组合,在二维平面上对应一个语料库 C
尚未占据的点。其设想是:把维度一推到最靠后以最大化
Δ\_save,同时用维度二的批量化与编码把 Δ\_recover 压到接近 certified
的水平,从而争取一个``低良好情形延迟且低恢复成本''的区域。DispersedLedger
已展示了编码在维度二上的威力,但其维度一定位与上述组合不同;Mysticeti/Mahi-Mahi
在维度一靠后,但维度二仍以``逐缺失触发''为主。上述组合点能否真被构造出一个兼具两种优势的协议,是本文坐标系提出的一个\textbf{可证伪的预测},留作未来工作------而能够提出这样一个具体、可证伪的预测,本身就是深层轴在``解释现有协议''之外、具有生成性的一面。

\subsection{8.5 深层轴的延伸:availability-enforcement
与排序公平性}\label{ux6df1ux5c42ux8f74ux7684ux5ef6ux4f38availability-enforcement-ux4e0eux6392ux5e8fux516cux5e73ux6027}

本文的深层轴虽以延迟为出发点,但它的影响并不止于延迟。DAG
共识的排序规则建立在 DAG 结构之上,而 DAG
结构由哪些区块可用、何时可用所塑造------这意味着
availability-enforcement
机制会间接影响交易的最终排序,进而触及最大可提取价值(MEV)与排序公平性问题 {[}53{]}--{[}56{]}。

本文的深层轴为理解这一关联提供了一个可检验的猜想:当 availability-enforcement
时机靠后时(post-commit),一个区块``是否进入
DAG''在更晚的时刻、由更依赖运行时条件的因素决定,这给了拜占庭节点更大的、操纵``哪些区块在排序时可见''的空间------即 enforcement 时机不仅是一个延迟维度,也是一个排序可操纵性维度。而在 pre-reference(certified) enforcement 下,区块内容在作者能观察到排序态势之前就已被认证锁定,这一操纵选项被关闭。须明确这只是一个待未来工作验证的机制猜想:它未量化敌手由此可获取的 MEV、未给出防护下界。但这一猜想展示了深层轴在延迟分析之外的潜在解释力------进一步支持了贡献 C 的可辩护内核:availability-enforcement 机制及其 reconciliation trigger,是比 certified/uncertified 更根本的变量。

\section{第 9 章
局限、误用声明与威胁有效性}\label{ux7b2c-9-ux7ae0-ux5c40ux9650ux8befux7528ux58f0ux660eux4e0eux5a01ux80c1ux6709ux6548ux6027}

\textbf{范围限定。} 本文第 5 章的归纳性结论限定在语料库
C(DAG-Rider、Narwhal/Tusk、Bullshark、Cordial
Miners、Mysticeti、Mahi-Mahi、Adelie 等)之内,不宣称对一切 DAG-BFT
协议普适。第 6 章命题 1
在纯异步模型下严格,部分同步下的推广本文仅作论证性讨论而未给出完整严格证明,这是一个明确的边界。

\textbf{实证仪器的两项已识别局限。}
本人在阶段七部署实证中识别出两项工具学限制,在此透明披露。其一,主动 RPC
探测器(rpc-probe)的高频轮询会与 fullnode 共识提交事件干涉,引入约 +30 ms
的系统性 cadence 偏差与显著方差;后续测量已切换到 fullnode
日志直接解析,不再依赖该探测器。其二,初版故障编排脚本对多个网络故障独立调用
\texttt{tc\ qdisc\ replace},在同一容器叠加多个 netem
故障时第二次调用会覆盖第一次,导致``delay + loss 同验证者''配置实际只产生
loss;该 bug 仅影响同容器叠加场景,§7.6.1
引用的跨验证者独立故障实证不受影响,修复后重跑的 Test 1 为完全可加(ratio
1.00)。

\textbf{独立种子复现性。} §7.6.1 的三个超可加候选 cell(Test 2、n=4
ablation、n=4 + 5\%)原始测量使用种子 20260525,本人以独立种子 20260526
完整复测(100 次受控重复,含 baseline/delay\_only/loss\_only/combo
完整对照链)。Test 2(n=7,5\% 丢包,1.35→1.08)与 n=4 +
5\%(1.17→1.12)超可加方向一致;而 n=4 ablation(2\% 丢包,故障数
\textgreater{} f)的原始超可加(1.16)未复现(复测
0.97,干净可加)------原结论为单种子假阳性。复测的关键修正是:超可加由故障强度驱动,而非故障数是否超越
BFT
边界。此修正本身即``本文实证结论需更多独立种子验证''这一局限的一个具体实例;其意图是为未来工作提供工具学与复现性诚信参考,而非为已有结论作事后辩护。

\textbf{注入式旋钮的构念效度。} §7.6.2 的 Δ\_recover
标定使用一个在共识恢复路径人为植入的时延旋钮来模拟单次恢复成本,并以轮转故障把恢复强制置于关键路径。两点须透明声明:其一,该旋钮是
emulation 而非对真实 certified 协议恢复开销的测量,故其标定的是或有负债沿
Δ\_recover 轴的形状(近似线性、委员会结构依赖),不可外推为 §7.2 中 Δ\_save
或 ρ*
的数值;其二,``恢复须处于关键路径方可观测''本身是一个实证发现(单节点恢复被推送式补块绕过),它精化而非削弱
§7.3 的或有负债论断。§7.6.3 进一步补上了一条 certified
对照臂(同引擎、可开关的
certificate-gate),其构念效度限定与本段同源并更强:该 gate
捕获的是认证的可用性证书(2f+1 接受,经 RoundProber
确认通道观测)而非密码学 2f+1 签名,故其 Δ\_save
是独立确认轮的延迟代价、其 ρ*
是成本结构意义上的翻转点,均不可外推为某真实 RBC 协议的绝对数值。

\textbf{条件化结论三条翻转轴的部署锚点及其共享限定。} 除 §7.6
的成本模型标定外,本文为第 7 章三条翻转条件与第 4--5
章的不变量分类各补了一项直接部署锚点:恢复频率轴(§7.3,cadence=3.78−265·ρ)、部分同步轴(§7.4,sub-critical
抖动被吸收的阴性结果)、网络非对称轴(§7.5,certified
臂更鲁棒)、逻辑型/物理型二分(§4.5)、以及 ①/②
分类(§5.2--5.4,逻辑型路径恒定在线 vs
物理型补偿按需爆发)。这些锚点共享一组须透明披露的限定:除 §7.5 复用了
§7.6.3 的可用性证书 gate(非密码学签名)作 certified 对照外,其余均在单一
uncertified 引擎(Sui/Mysticeti)上展开、无 certified
对照臂,且涉及恢复成本处皆使用注入式 Δ\_recover 旋钮;§4.5
的物理型一半受协议块大小限制,字节量效应在可测区间(≤10 KB)内未能与 RPC
数量分离;§7.4 为诚实阴性结果(sub-critical
设计下信号本就弱)。因此这些锚点确证的是各轴上 uncertified
一侧成本形态的方向与量级,而非跨两套独立协议引擎的绝对延迟比较------后者仍属未来工作。

\textbf{模型抽象的代价。} 第 7 章的成本模型(净优势 ≈ Δ\_save −
ρ·Δ\_recover)是一个有意简化的定性模型。它假定补偿成本可加性地计入,忽略了恢复请求之间可能的相互干扰与排队效应。这一简化使模型易于分析,但也意味着它给出的是趋势而非精确预测;真实系统中
ρ 与 Δ\_recover 还可能相互耦合(如 7.4 节的拥塞反馈)。

\textbf{framing 风险。} 贡献 C 是一次 framing
工程:它把一个隐式贯穿全文的深层轴提升为显式坐标。这类工作的失败模式不是``定理塌了'',而是``读者没接住''。本文通过把深层轴操作化为两个可测维度来降低这一风险,但无法完全消除。

\textbf{误用风险。} 如第 1.4
节所述,本文的条件化结论存在被断章引用的风险。本文重申:任何``uncertified
是否更优''的判断都不应脱离其网络与敌手模型前提。相比之下,贡献 C
由于陈述的是一个设计轴而非一个比较句,在结构上更不易被断章------这一性质本身也是本文的一个附带观察。

\textbf{未充分展开的利益相关者。}
本文的成本台账以验证者为中心。两个利益相关者的视角值得未来工作纳入:其一,轻客户端与外部验证者------数据可用性对无法存储全量数据的轻客户端尤为关键,availability-enforcement
时机靠后意味着轻客户端更晚才能确信数据可得,这与数据可用性采样及欺诈/可用性证明
{[}48{]} 的研究直接相关;其二,验证者运营经济学------reconciliation
的恢复带宽由谁付费、如何计价,是真实部署决策中的现实变量,本文的定性模型未予刻画。把这两类视角纳入成本台账,会使第
7 章的条件化结论更贴近部署现实。

\textbf{威胁有效性。}
本文的核心威胁来自成本台账的参数化:Δ\_save、Δ\_recover、ρ
的具体取值依赖于协议实现与部署环境,本文给出的是其定性关系而非定量测量。本文已对其中的或有负债侧提供了部署探查(§7.6.1、§7.6.2),并以一个同引擎、可开关的
certificate-gate 构建了 certified 对照臂(§7.6.3),首次在部署上测得正的
Δ\_save、认证对恢复成本的缩减、以及净优势翻转的 ρ*
band。须明确这条对照臂的限定:它实现的是认证的可用性证书(2f+1
接受)而非密码学签名聚合,仍为单引擎、Δ\_recover
为注入式旋钮,故标定的是净优势的成本结构而非真实协议对的绝对延迟。一个以真实
RBC 签名开销为对照、跨两套独立引擎的受控 benchmark
仍属未来工作。本文的论证也依赖对若干协议安全性证明的解读,解读虽力求忠实,但不排除存在偏差;本文已尽量在引用处标注所依赖的具体性质。

\section{第 10 章
结论与未来工作}\label{ux7b2c-10-ux7ae0-ux7ed3ux8bbaux4e0eux672aux6765ux5de5ux4f5c}

本文对 DAG-based BFT 共识中一个被广泛接受的判断------``uncertified DAG
通过去认证免费降低延迟''------提出了系统性的质疑。

回答主研究问题:uncertified DAG
的低延迟优势成立于一个明确界定的条件区域(网络近同步、故障稀少、负载均衡,即在
§7.2 独立性假设下恢复频率 ρ 低于临界值 ρ*
的区域),在该区域之外会被隐性代价抵消甚至反转。回答
SubRQ1:认证维持非等价性、数据可用性、因果历史完整性三个不变量,其中可用性是唯一的物理型不变量。回答
SubRQ2:去认证后,逻辑型不变量被低成本平移,物理型的数据可用性被弱化并补偿,隐性代价集中于此。回答
SubRQ3:延迟优势是条件性的,而非结构性的;部署实测进一步显示,该条件化的或有负债侧确为一笔非线性、由故障强度驱动的尾部代价(§7.6.1);本文进一步为三条翻转条件各补了一项直接部署锚点(恢复频率
§7.3、部分同步 §7.4、网络非对称 §7.5),连同 §4.5 的不变量异质性与 §5 的
①/②
分类实证,使条件化结论的每一维都获得部署可观测的支撑(其共享的单引擎/可用性证书/注入旋钮限定见
§9)。

本文的根本结论是:移除认证不是做减法,而是一次结构性替换;而
certified/uncertified 只是浅层标签,真正的深层设计轴是
availability-enforcement 机制与 reconciliation trigger
条件。本文进一步论证,这一深层轴中真正不可被归为重命名的内核------reconciliation
trigger
维度,以及延迟与鲁棒性的二维联合优化------不仅解释延迟,也触及排序公平性,并能提出可证伪的新设计点预测,因而具有超出原初问题的生成性与解释力。

对协议设计者而言,这意味着设计选择不应停留在``要不要认证'',而应直接面对深层轴的两个维度------根据目标部署环境的故障率与网络对称性,选择
availability-enforcement 的时机与 reconciliation 的触发策略。对
benchmark
与评测设计者而言,这意味着评测必须扩大采样,覆盖故障情形与非对称网络,否则报告的延迟数字只是高维性能函数的一个低维投影。

未来工作有两个方向:其一,把第 8
章的深层轴形式化,给出两个维度上的可计算度量,特别是 reconciliation
trigger 条件的形式刻画;其二,把第 7
章成本台账提升为定量的区域划分。本文已在单一 uncertified 协议(Sui
n=7/n=4)上对或有负债侧作了首个部署探查(§7.6.1、§7.6.2):§7.6.1
证实该代价为非线性、由故障强度驱动,§7.6.2
进一步以受控旋钮标定出或有负债沿 Δ\_recover
轴近似线性、随委员会收缩而加剧、且仅在恢复处于关键路径时
materialize;§7.6.3 则以一个同引擎、可开关的 certificate-gate 补上
certified 对照臂,首次在部署上把 Δ\_save
测为正、确认认证缩减恢复成本、并把净优势翻转 ρ* 定量定位于约
0.013--0.038 次/commit。尚待完成的是把这条可用性证书的对照臂升级为含真实
RBC 签名开销、跨两套独立引擎的受控
benchmark,以从``成本结构标定''推进到``真实协议对的绝对延迟比较''。

\section{声明}\label{sec:declaration}

\textbf{数据可用性声明(Data Availability
Statement)}:本文以理论解析为内核,并含一个部署实证组件。文中引用的全部协议规范与论文均为公开发表文献,可经由参考文献所列
DOI、arXiv 或 IACR ePrint 编号获取。本文全部代码、数据与实验材料已公开存档于
GitHub 仓库
\texttt{https://\allowbreak{}github.com/\allowbreak{}Lang-Qiu/\allowbreak{}certified-vs-uncertified-dag-bft},以下按模块说明。

\textbf{仿真代码与解析模型}:第 7--8
章成本模型的仿真与图表生成代码(Python/SimPy,含成本模型、Monte Carlo 仿真与
Figure 1--3 的可复现脚本)存档于仓库
\texttt{submit/\allowbreak{}simulation/\allowbreak{}};图 1、图 2
为该解析模型的数值图示,图 3 为按第 8
章深层轴框架的定性定位,均非真实协议实测。形式化模型代码存档于
\texttt{submit/\allowbreak{}model/\allowbreak{}dag\_\allowbreak{}bft\_\allowbreak{}model.py}。

\textbf{部署实证}:§7.6.1 的实证(图 7-3、表
7-1)来自真实 Sui n=7/n=4 多验证者部署的受控故障注入实验,共 290+
次重复(含独立种子复现)。部署基于 MystenLabs/sui 主分支 commit
\texttt{62ee6ada}(2026-05-22)。实验编排脚本、Docker Compose 配置与分析工具存档于仓库
\texttt{submit/\allowbreak{}deployment/\allowbreak{}multivalidator/\allowbreak{}};503 次逐 run
指标(含 \texttt{consensus\_metrics.csv}、控制器记录与配置快照)存档于
\texttt{submit/\allowbreak{}deployment/\allowbreak{}multivalidator/\allowbreak{}data/\allowbreak{}per\_\allowbreak{}run\_\allowbreak{}metrics/\allowbreak{}};19
个汇总 CSV 存档于
\texttt{submit/\allowbreak{}deployment/\allowbreak{}multivalidator/\allowbreak{}data/\allowbreak{}runs/\allowbreak{}}。

§7.6.2 的恢复侧旋钮标定(图
7-4)来自带受控 $\Delta$\_recover 旋钮的 Sui
构建在关键路径轮转故障下的部署实测,其 harness 与图表生成脚本存档于
\texttt{submit/\allowbreak{}calibration/\allowbreak{}scripts/\allowbreak{}}。§7.6.3 的
certified 对照臂标定(图 7-5)来自同一 Sui 构建中一个可经环境变量开关的
certificate-gate(Route
A1,认证的可用性证书操作化)在关键路径轮转故障下的部署实测,其共识源码改动、Phase
1/2 harness 与图表脚本亦存档于
\texttt{submit/\allowbreak{}calibration/\allowbreak{}}。

此外,第 4--5 章与
§7.3--7.5 的部署锚点亦来自真实 Sui n=7/n=4 多验证者部署的受控实验:§4.5
区块负载异质性(图 4-1)、§5.2--5.3 逻辑型路径固定开销(图 5-1)、§5.4
数据可用性或有负债的丢包/分区对照与剂量响应(图 5-2、图 5-3)、§7.3
cadence 随恢复频率退化(图 7-6)、§7.4 sub-critical 抖动阴性结果(图
7-7)、§7.5 非对称网络下 certified 臂鲁棒性(图 7-8),其编排脚本、逐 run
汇总 CSV(\texttt{payload\_size\_summary.csv}、\texttt{chapter5\_enforcement\_summary.csv}、\texttt{partition\_grid\_summary.csv}、\texttt{rho\_spacing\_summary.csv}、\texttt{rho\_scan\_summary.csv}、\texttt{gst\_jitter\_summary.csv}、\texttt{net\_asymmetry\_summary.csv})、分析与图表生成脚本均存档于仓库
\texttt{submit/\allowbreak{}calibration/\allowbreak{}} 与
\texttt{submit/\allowbreak{}deployment/\allowbreak{}}。

两起数据完整性事件(P0-1 首轮
\texttt{timeout.exe} 伪迹致 recFetch 误记为 0、P0-2 $\rho$ 首轮 $\rho$
近恒定的设计缺陷)的诚实记录与受污染数据存档见仓库
\texttt{submit/\allowbreak{}calibration/\allowbreak{}EXPERIMENTAL\_\allowbreak{}COVERAGE.md}。

\textbf{作者贡献(CRediT)}:概念化、方法论、形式化分析、写作------原稿与修订,均由论文作者完成。

\textbf{AI 使用声明(AI
Disclosure)}:本文的结构组织和实验脚本撰写使用了大语言模型辅助。研究论点、核心论证(包括第 6 章必要性论证与第 8
章深层轴的提出)与最终判断由作者负责并确认。全部参考文献经独立检索核验。

\begin{thebibliography}{99}
\setlength{\itemsep}{2pt plus 1pt}
\bibitem{ref1} L. Lamport, R. Shostak, and M. Pease, “The Byzantine generals problem,” \textit{ACM Trans. Program. Lang. Syst.}, vol. 4, no. 3, pp. 382–401, 1982.
\bibitem{ref2} M. J. Fischer, N. A. Lynch, and M. S. Paterson, “Impossibility of distributed consensus with one faulty process,” \textit{J. ACM}, vol. 32, no. 2, pp. 374–382, 1985.
\bibitem{ref3} C. Dwork, N. Lynch, and L. Stockmeyer, “Consensus in the presence of partial synchrony,” \textit{J. ACM}, vol. 35, no. 2, pp. 288–323, 1988.
\bibitem{ref4} F. B. Schneider, “Implementing fault-tolerant services using the state machine approach: A tutorial,” \textit{ACM Comput. Surv.}, vol. 22, no. 4, pp. 299–319, 1990.
\bibitem{ref5} G. Bracha, “Asynchronous Byzantine agreement protocols,” \textit{Inf. Comput.}, vol. 75, no. 2, pp. 130–143, 1987.
\bibitem{ref6} C. Cachin and S. Tessaro, “Asynchronous verifiable information dispersal,” in \textit{Proc. 24th IEEE Symp. Reliable Distributed Systems (SRDS)}, 2005, pp. 191–202.
\bibitem{ref7} C. Cachin, K. Kursawe, and V. Shoup, “Random oracles in Constantinople: Practical asynchronous Byzantine agreement using cryptography,” \textit{J. Cryptol.}, vol. 18, no. 3, 2005.
\bibitem{ref8} D. Boneh, B. Lynn, and H. Shacham, “Short signatures from the Weil pairing,” in \textit{Proc. ASIACRYPT}, 2001.
\bibitem{ref9} C. Cachin, R. Guerraoui, and L. Rodrigues, \textit{Introduction to Reliable and Secure Distributed Programming}, 2nd ed. Berlin, Germany: Springer, 2011.
\bibitem{ref10} M. Castro and B. Liskov, “Practical Byzantine fault tolerance,” in \textit{Proc. 3rd USENIX Symp. Operating Systems Design and Implementation (OSDI)}, 1999.
\bibitem{ref11} M. Yin, D. Malkhi, M. K. Reiter, G. Golan-Gueta, and I. Abraham, “HotStuff: BFT consensus with linearity and responsiveness,” in \textit{Proc. ACM Symp. Principles of Distributed Computing (PODC)}, 2019, pp. 347–356.
\bibitem{ref12} I. Abraham, D. Malkhi, K. Nayak, L. Ren, and M. Yin, “Sync HotStuff: Simple and practical synchronous state machine replication,” in \textit{Proc. IEEE Symp. Security and Privacy (S\&P)}, 2020.
\bibitem{ref13} E. Buchman, J. Kwon, and Z. Milosevic, “The latest gossip on BFT consensus,” arXiv:1807.04938, 2018.
\bibitem{ref14} C. Stathakopoulou, T. David, M. Pavlovic, and M. Vukolić, “Mir-BFT: High-throughput robust BFT for decentralized networks,” arXiv:1906.05552, 2019.
\bibitem{ref15} J. Camenisch, M. Drijvers, T. Hanke, Y.-A. Pignolet, V. Shoup, and D. Williams, “Internet Computer Consensus,” in \textit{Proc. ACM Symp. Principles of Distributed Computing (PODC)}, 2022, pp. 81–91.
\bibitem{ref16} T. Crain, V. Gramoli, M. Larrea, and M. Raynal, “DBFT: Efficient leaderless Byzantine consensus and its application to blockchains,” in \textit{Proc. 17th IEEE Int. Symp. Network Computing and Applications (NCA)}, 2018.
\bibitem{ref17} A. Miller, Y. Xia, K. Croman, E. Shi, and D. Song, “The honey badger of BFT protocols,” in \textit{Proc. ACM SIGSAC Conf. Computer and Communications Security (CCS)}, 2016.
\bibitem{ref18} S. Duan, M. K. Reiter, and H. Zhang, “BEAT: Asynchronous BFT made practical,” in \textit{Proc. ACM SIGSAC Conf. Computer and Communications Security (CCS)}, 2018, pp. 2028–2041.
\bibitem{ref19} B. Guo, Z. Lu, Q. Tang, J. Xu, and Z. Zhang, “Dumbo: Faster asynchronous BFT protocols,” in \textit{Proc. ACM SIGSAC Conf. Computer and Communications Security (CCS)}, 2020, pp. 803–818.
\bibitem{ref20} I. Abraham, D. Malkhi, and A. Spiegelman, “Asymptotically optimal validated asynchronous Byzantine agreement,” in \textit{Proc. ACM Symp. Principles of Distributed Computing (PODC)}, 2019.
\bibitem{ref21} Y. Lu, Z. Lu, Q. Tang, and G. Wang, “Dumbo-MVBA: Optimal multi-valued validated asynchronous Byzantine agreement, revisited,” in \textit{Proc. ACM Symp. Principles of Distributed Computing (PODC)}, 2020.
\bibitem{ref22} S. Nakamoto, “Bitcoin: A peer-to-peer electronic cash system,” white paper, 2008.
\bibitem{ref23} J. A. Garay, A. Kiayias, and N. Leonardos, “The Bitcoin backbone protocol: Analysis and applications,” in \textit{Proc. EUROCRYPT}, 2015, pp. 281–310.
\bibitem{ref24} Y. Sompolinsky and A. Zohar, “Secure high-rate transaction processing in Bitcoin,” in \textit{Proc. Financial Cryptography and Data Security (FC)}, 2015, pp. 507–527.
\bibitem{ref25} Y. Gilad, R. Hemo, S. Micali, G. Vlachos, and N. Zeldovich, “Algorand: Scaling Byzantine agreements for cryptocurrencies,” in \textit{Proc. 26th ACM Symp. Operating Systems Principles (SOSP)}, 2017, pp. 51–68.
\bibitem{ref26} Team Rocket, M. Yin, K. Sekniqi, R. van Renesse, and E. G. Sirer, “Scalable and probabilistic leaderless BFT consensus through metastability,” arXiv:1906.08936, 2019.
\bibitem{ref27} V. Bagaria, S. Kannan, D. Tse, G. Fanti, and P. Viswanath, “Prism: Deconstructing the blockchain to approach physical limits,” in \textit{Proc. ACM SIGSAC Conf. Computer and Communications Security (CCS)}, 2019.
\bibitem{ref28} J. Neu, E. N. Tas, and D. Tse, “Ebb-and-flow protocols: A resolution of the availability-finality dilemma,” in \textit{Proc. IEEE Symp. Security and Privacy (S\&P)}, 2021.
\bibitem{ref29} V. Buterin, “Ethereum: A next-generation smart contract and decentralized application platform,” white paper, 2014.
\bibitem{ref30} V. Buterin and V. Griffith, “Casper the friendly finality gadget,” arXiv:1710.09437, 2017.
\bibitem{ref31} I. Keidar, E. Kokoris-Kogias, O. Naor, and A. Spiegelman, “All you need is DAG,” in \textit{Proc. ACM Symp. Principles of Distributed Computing (PODC)}, 2021.
\bibitem{ref32} G. Danezis, L. Kokoris-Kogias, A. Sonnino, and A. Spiegelman, “Narwhal and Tusk: A DAG-based mempool and efficient BFT consensus,” in \textit{Proc. 17th European Conf. Computer Systems (EuroSys)}, 2022.
\bibitem{ref33} A. Spiegelman, N. Giridharan, A. Sonnino, and L. Kokoris-Kogias, “Bullshark: DAG BFT protocols made practical,” in \textit{Proc. ACM SIGSAC Conf. Computer and Communications Security (CCS)}, 2022.
\bibitem{ref34} A. Gągol, D. Leśniak, D. Straszak, and M. Świętek, “Aleph: Efficient atomic broadcast in asynchronous networks with Byzantine nodes,” in \textit{Proc. 1st ACM Conf. Advances in Financial Technologies (AFT)}, 2019.
\bibitem{ref35} X. Dai, Z. Zhang, J. Xiao, J. Yue, X. Xie, and H. Jin, “GradedDAG: An asynchronous DAG-based BFT consensus with lower latency,” in \textit{Proc. 42nd Int. Symp. Reliable Distributed Systems (SRDS)}, 2023.
\bibitem{ref36} X. Dai, G. Wang, J. Xiao, Z. Guo, R. Hao, X. Xie, and H. Jin, “LightDAG: A low-latency DAG-based BFT consensus through lightweight broadcast,” in \textit{Proc. IEEE Int. Parallel and Distributed Processing Symp. (IPDPS)}, 2024, pp. 998–1008.
\bibitem{ref37} A. Spiegelman, B. Arun, R. Gelashvili, and Z. Li, “Shoal: Improving DAG-BFT latency and robustness,” in \textit{Proc. Financial Cryptography and Data Security (FC)}, 2024.
\bibitem{ref38} B. Arun, Z. Li, F. Suri-Payer, S. Das, and A. Spiegelman, “Shoal++: High throughput DAG BFT can be fast!,” in \textit{Proc. 22nd USENIX Symp. Networked Systems Design and Implementation (NSDI)}, 2025.
\bibitem{ref39} N. Shrestha, R. Shrothrium, A. Kate, and K. Nayak, “Sailfish: Towards improving the latency of DAG-based BFT,” in \textit{Proc. IEEE Symp. Security and Privacy (S\&P)}, 2025.
\bibitem{ref40} D. Malkhi, C. Stathakopoulou, and M. Yin, “BBCA-Chain: Low latency, high throughput BFT consensus on a DAG,” in \textit{Proc. Financial Cryptography and Data Security (FC)}, 2024.
\bibitem{ref41} N. Shrestha and A. Kate, “Towards improving throughput and scalability of DAG-based BFT SMR,” IACR Cryptol. ePrint Arch., Paper 2025/877, 2025.
\bibitem{ref42} I. Keidar, O. Naor, O. Poupko, and E. Shapiro, “Cordial Miners: Fast and efficient consensus for every eventuality,” in \textit{Proc. 37th Int. Symp. Distributed Computing (DISC)}, 2023.
\bibitem{ref43} K. Babel, A. Chursin, G. Danezis, L. Kokoris-Kogias, and A. Sonnino, “Mysticeti: Reaching the limits of latency with uncertified DAGs,” in \textit{Proc. Network and Distributed System Security Symp. (NDSS)}, 2025.
\bibitem{ref44} P. Jovanovic, L. Kokoris-Kogias, B. Kumara, A. Sonnino, P. Tennage, and I. Zablotski, “Mahi-Mahi: Low-latency asynchronous BFT DAG-based consensus,” in \textit{Proc. IEEE Int. Conf. Distributed Computing Systems (ICDCS)}, 2025, pp. 549–559.
\bibitem{ref45} A. Chursin, “Adelie: Detection and prevention of Byzantine behaviour in DAG-based consensus protocols,” arXiv:2408.02000, 2024.
\bibitem{ref46} S. Blackshear, A. Chursin, G. Danezis, A. Kichidis, L. Kokoris-Kogias, X. Li, M. Logan, et al., “Sui Lutris: A blockchain combining broadcast and consensus,” in \textit{Proc. ACM SIGSAC Conf. Computer and Communications Security (CCS)}, 2024.
\bibitem{ref47} L. Yang, S. J. Park, M. Alizadeh, S. Kannan, and D. Tse, “DispersedLedger: High-throughput Byzantine consensus on variable bandwidth networks,” in \textit{Proc. 19th USENIX Symp. Networked Systems Design and Implementation (NSDI)}, 2022.
\bibitem{ref48} M. Al-Bassam, A. Sonnino, and V. Buterin, “Fraud and data availability proofs: Detecting invalid blocks in light clients,” in \textit{Proc. Financial Cryptography and Data Security (FC)}, 2021.
\bibitem{ref49} N. Giridharan, F. Suri-Payer, I. Abraham, L. Alvisi, and N. Crooks, “Autobahn: Seamless high speed BFT,” in \textit{Proc. 30th ACM Symp. Operating Systems Principles (SOSP)}, 2024.
\bibitem{ref50} E. Kokoris-Kogias, P. Jovanovic, L. Gasser, N. Gailly, E. Syta, and B. Ford, “OmniLedger: A secure, scale-out, decentralized ledger via sharding,” in \textit{Proc. IEEE Symp. Security and Privacy (S\&P)}, 2018, pp. 583–598.
\bibitem{ref51} P. Civit, S. Gilbert, and V. Gramoli, “Polygraph: Accountable Byzantine agreement,” in \textit{Proc. IEEE Int. Conf. Distributed Computing Systems (ICDCS)}, 2021.
\bibitem{ref52} P. Civit, S. Gilbert, V. Gramoli, R. Guerraoui, and J. Komatovic, “As easy as ABC: Optimal (A)ccountable (B)yzantine (C)onsensus is easy!,” in \textit{Proc. IEEE Int. Parallel and Distributed Processing Symp. (IPDPS)}, 2022.
\bibitem{ref53} P. Daian, S. Goldfeder, T. Kell, Y. Li, X. Zhao, I. Bentov, L. Breidenbach, and A. Juels, “Flash Boys 2.0: Frontrunning, transaction reordering, and consensus instability in decentralized exchanges,” in \textit{Proc. IEEE Symp. Security and Privacy (S\&P)}, 2020.
\bibitem{ref54} M. Kelkar, F. Zhang, S. Goldfeder, and A. Juels, “Order-fairness for Byzantine consensus,” in \textit{Proc. CRYPTO}, 2020.
\bibitem{ref55} M. Kelkar, S. Deb, S. Long, A. Juels, and S. Kannan, “Themis: Fast, strong order-fairness in Byzantine consensus,” in \textit{Proc. ACM SIGSAC Conf. Computer and Communications Security (CCS)}, 2023.
\bibitem{ref56} D. Kang, J. Chen, T. T. A. Dinh, and M. Sadoghi, “FairDAG: Consensus fairness over multi-proposer causal design,” \textit{Proc. VLDB Endowment}, vol. 19, 2026.
\bibitem{ref57} M. Raikwar, N. Polyanskii, and S. Müller, “SoK: DAG-based consensus protocols,” in \textit{Proc. IEEE Int. Conf. Blockchain and Cryptocurrency (ICBC)}, 2024.
\bibitem{ref58} S. Bano, A. Sonnino, M. Al-Bassam, S. Azouvi, P. McCorry, S. Meiklejohn, and G. Danezis, “SoK: Consensus in the age of blockchains,” in \textit{Proc. 1st ACM Conf. Advances in Financial Technologies (AFT)}, 2019.
\bibitem{ref59} J. A. Garay and A. Kiayias, “SoK: A consensus taxonomy in the blockchain era,” in \textit{Proc. CT-RSA}, 2020.
\bibitem{ref60} C. Cachin and M. Vukolić, “Blockchain consensus protocols in the wild,” in \textit{Proc. 31st Int. Symp. Distributed Computing (DISC)}, 2017.
\bibitem{ref61} L. Baird, “The Swirlds hashgraph consensus algorithm: Fair, fast, Byzantine fault tolerance,” Swirlds, Tech. Rep. SWIRLDS-TR-2016-01, 2016.
\bibitem{ref62} T. Crain, C. Natoli, and V. Gramoli, “Red Belly: A secure, fair and scalable open blockchain,” in \textit{Proc. IEEE Symp. Security and Privacy (S\&P)}, 2021, pp. 466–483.
\end{thebibliography}



\end{document}
